%
%  PARA TRABALLOS EN GALLEGO USAR (LINEA 12): \usepackage[galician]{babel}
%  PARA TRABALLOS EN CASTELLANO USAR (LINEA 13): \usepackage[spanish]{babel}
%
% Para los acentos usamos codificacion UTF-8 (LINEA 10): \usepackage[utf8]{inputenc} 
% Si se usase la codificacion es_ES.ISO-8859-1 (LINEA 11): \usepackage[latin1]{inputenc}
% La conversion de acentos se hace con: iconv -f UTF-8 -t ISO-8859-1 filename.tex
%
% Como se incluyen figuras eps hay que compilar con: latex traballo , dvipdf traballo
%

\documentclass[12pt,twoside,a4paper]{book}
% pódense engadir todos os packages necesarios
\usepackage[utf8]{inputenc}
% \usepackage[latin1]{inputenc}
\usepackage[spanish]{babel}
\usepackage{graphicx}
\usepackage[dvips]{epsfig}
\usepackage{amssymb}
\usepackage{eurosym}
\usepackage{float}
\usepackage{latexsym}
\usepackage{a4}
\usepackage{listings}
\usepackage{hyperref} % menús no pdf pero non leva ben co package galician


\addto\captionsgalician{\def\contentsname{Memoria tipo B -- \'{I}ndice xeral }}
	

\begin{document}
\pagestyle{empty}
\begin{center}
	{\bf\Large UNIVERSIDADE DE SANTIAGO DE COMPOSTELA}
	
	\vspace{0.5cm}
	\includegraphics[width=5cm]{figuras/logo_usc.eps}
	
	\vspace{0.5cm}
	{\bf\large ESCOLA TÉCNICA SUPERIOR DE ENXEÑARÍA}
	
	\vspace{3cm}
	{\bf\LARGE Herramienta para la creación y edición colaborativa de documentos}
	
	%\vspace{0.5cm}
	%{\bf\LARGE Subtítulo do Traballo de Fin de Grao}
\end{center}

\vspace{2cm}
\hspace{4cm}\begin{tabular}{l}
	{\it\Large Autor:} \\
	{\bf\Large Pablo Tarrío Otero} \\
	~ \\
	{\it\Large Tutores:} \\
	{\bf\Large Manuel Lama Penín} \\
	{\bf\Large Juan Carlos Vidal Aguiar} \\
	{\bf\Large Víctor José Gallego Fontenla} \\
\end{tabular}

\vspace{2cm}
\begin{center}
	{\bf\Large Grado en Ingeniería Informática}
	
	\vspace{0.5cm}
	{\bf\large Julio 2022}
	
	\vspace{0.5cm}
	Trabajo de Fin de Grado presentado en la {\it Escola Técnica Superior de Enxeñaría} de la {\it Universidade de Santiago de Compostela} para la obtención del Grado en Ingeniería Informática
\end{center}
\cleardoublepage
\pagestyle{plain}
\pagenumbering{roman}
\includegraphics[width=4cm]{figuras/logo_usc.eps}

\vspace{1cm}
{\bf D. Manuel Lama Penín}, Profesor del Departamento de Electrónica y Computación de la {\it Universidade de Santiago de Compostela}, {\bf D. Juan Carlos Vidal Aguiar}, Profesor del Departamento de Electrónica y Computación de la {\it Universidade de Santiago de Compostela}, y {\bf D. Víctor José Gallego Fontenla}, Investigador Predoctoral FPU del {\it Ministerio de Ciencia, Innovación y Universidades},

\vspace{1cm}
INFORMAN:

\vspace{1cm}
Que la presente memoria, titulada {\it Herramienta para la creación y edición
colaborativa de documentos}, presentada por {\bf D. Pablo Tarrío Otero} para superar los créditos correspondientes al Trabajo de Fin de Grado de la titulación de Grado en Ingeniería Informática, se realizó bajo nuestra tutoría en el Departamento de Electrónica y Computación de la {\it Universidade de Santiago de Compostela}.

\vspace{1cm}
Y para que así conste a los efectos oportunos, expiden el presente informe en Santiago de Compostela, a 21 de julio de 2022:

\vspace{0.5cm}
\begin{tabular}{ll}
	Tutor, & \hspace{3.6cm} Cotutor, \\
	~ \\
	~ \\
	~ \\
	~ \\
	Manuel Lama Penín & \hspace{3.6cm} Juan Carlos Vidal Aguiar
\end{tabular}

\vspace{1cm}
\begin{tabular}{ll}
	Cotutor, & \hspace{2cm} Alumno, \\
	~ \\
	~ \\
	~ \\
	~ \\
	Víctor José Gallego Fontenla & \hspace{2cm} Pablo Tarrío Otero
\end{tabular}

 % paxina de certificación (optativa)
\cleardoublepage
\pagestyle{plain}
\chapter*{Agradecimientos}

{\it 
A mis tutores Manuel, Juan Carlos y Víctor por haberme orientado durante estos meses a la hora de realizar el TFG. Sin ellos el resultado de este proyecto no habría sido tan satisfactorio para mí.
\\

A mis amigos, por haberme acompañado durante todos estos años y haberme hecho pasar una etapa preciosa, que aunque ha sido corta, resultará inolvidable en mi vida.
\\

A mi familia, por no dejar que me rindiese ante cualquier bache que se presentase, y por haberme animado durante toda mi vida desde que era un niño hasta el día de hoy.
} % paxina de agradecementos (optativa) 
\cleardoublepage
\pagestyle{plain}
\chapter*{Resumen}
Este trabajo de fin de grado aborda el desarrollo de una herramienta colaborativa que permitirá la creación, edición y consolidación de documentos por parte de varios usuarios. Esta herramienta será integrada en una plataforma de consulta de legislación consolidada gallega, facilitando la importación de legislación desde el Diario Oficial de Galicia (DOG) \cite{dog}, así como el mantenimiento con la aplicación en las sucesivas modificaciones que se vayan realizando.
\\

De esta forma, los usuarios podrán importar documentos a la herramienta colaborativa a partir de leyes importadas del DOG. Por una parte, se permite la búsqueda de leyes y normas encontradas en el DOG filtrando según los parámetros deseados. Adicionalmente, podrá verse una previsualización en la página oficial del DOG. Dichos documentos encontrados en el DOG podrán ser importados a la herramienta colaborativa, donde se podrá proceder a su edición. En cuanto a la edición, se ofrece la posibilidad de proponer cambios a la norma que se va a editar, manipular las leyes vinculadas a esta, así como escribir posibles notas aclaratorias. Además, todo ello será gestionado por usuarios previamente registrados. % páxina de resumo (optativa) 

\cleardoublepage
\pagestyle{plain}
\tableofcontents
\listoffigures
\listoftables

% Agora incluimos os capítulos. Cambiamos a numeración e as cabeceiras
\cleardoublepage
\pagenumbering{arabic}
\setcounter{page}{1}
\pagestyle{headings}
\chapter{Especificación de requisitos}
\label{enlaceespecificacion}

Para proceder a especificar los requisitos se ha realizado un profundo análisis de los objetivos de la aplicación. Estos requisitos se han buscado teniendo en cuenta aspectos funcionales, de rendimiento, de seguridad, de información, etc.
\\

Las plantillas empleadas para las tablas de los requisitos y casos de uso se encuentran en los apéndices: la \hyperref[enlaceFRQX]{E.1. Plantilla de requisitos funcionales}, la \hyperref[enlaceNFRX]{E.2. Plantilla de requisitos no funcionales}, la \hyperref[enlaceIRQX]{E.3. Plantilla de requisitos de información} y la \hyperref[enlaceUCX]{E.4. Plantilla de casos de uso}.
Estas plantillas están inspiradas en las tablas empleadas en el REM \cite{rem} para la gestión de requisitos.
\\

La importancia de los requisitos será medida según el grado de la necesidad del requisito en la aplicación, siendo en el caso de este proyecto ``Vital" para los requisitos de mayor importancia, y, ``Quedaría bien" para los menos importantes. Por su parte, la urgencia, indicará aquellas partes que se han de implementar con mayor rapidez, siendo ``Inmediatamente" las que tienen una mayor prioridad, y, ``Puede esperar", para los requisitos que no son estrictamente necesarios.

\section{Requisitos funcionales}

Los requisitos funcionales \cite{requisitos} son aquellos que definen una funcionalidad del sistema de software o sus componentes, tanto para las entradas de datos como para el comportamiento en situaciones específicas.
\\

En el caso de esta aplicación se han localizado dos grupos de requisitos funcionales: los asociados con operaciones de usuarios, y los que se relacionan con operaciones asociadas a leyes.
\\

Los requisitos funcionales pueden ser consultados en el \hyperref[APRequisitosFuncionales]{Apéndice C.1. Requisitos funcionales}.
\section{Requisitos no funcionales}
\label{APRequisitosNoFuncionales}

\begin{table}[H]
\begin{center}
\begin{tabular}{|p{3cm}|p{10cm}|} \hline
\centering {\bf NFR-01} & Tiempo de respuesta $<$ 2 segundos  \\ \hline\hline
\centering {\bf Versión} & 1.0 (05/06/2022) \\ \hline
\centering {\bf Descripción} & El sistema deberá responder en menos de 2 segundos, sin tener en cuenta el tiempo que tarde la red en enviar la información de vuelta al usuario. \\ \hline
\centering {\bf Importancia} & Vital \\ \hline
\centering {\bf Urgencia} & Inmediatamente \\ \hline
\end{tabular}
\caption{Requisito no funcional 01. Tiempo de respuesta $<$ 2 segundos.}
\label{enlaceNFR1}
\end{center}
\end{table}

\begin{table}[H]
\begin{center}
\begin{tabular}{|p{3cm}|p{10cm}|} \hline
\centering {\bf NFR-02} & Los datos sensibles del usuario se almacenan encriptados  \\ \hline\hline
\centering {\bf Versión} & 1.0 (05/06/2022) \\ \hline
\centering {\bf Descripción} & El sistema deberá encriptar los datos sensibles del usuario como las contraseñas en la base de datos de lex.gal. Se utilizará un algoritmo SHA-256 para encriptar. \\ \hline
\centering {\bf Importancia} & Vital \\ \hline
\centering {\bf Urgencia} & Inmediatamente \\ \hline
\end{tabular}
\caption{Requisito no funcional 02. Los datos sensibles del usuario se almacenan encriptados.}
\label{enlaceNFR2}
\end{center}
\end{table}

\begin{table}[H]
\begin{center}
\begin{tabular}{|p{3cm}|p{10cm}|} \hline
\centering {\bf NFR-03} & No permitir que un rol realice operaciones de otro  \\ \hline\hline
\centering {\bf Versión} & 1.0 (05/06/2022) \\ \hline
\centering {\bf Descripción} & El sistema deberá restringir aquellas operaciones exclusivas del administrador para el resto de usuarios. \\ \hline
\centering {\bf Importancia} & Vital \\ \hline
\centering {\bf Urgencia} & Inmediatamente \\ \hline
\end{tabular}
\caption{Requisito no funcional 03. No permitir que un rol realice operaciones de otro.}
\label{enlaceNFR3}
\end{center}
\end{table}

\begin{table}[H]
\begin{center}
\begin{tabular}{|p{3cm}|p{10cm}|} \hline
\centering {\bf NFR-04} & Permanecer operativo durante el horario de trabajo \\ \hline\hline
\centering {\bf Versión} & 1.0 (05/06/2022) \\ \hline
\centering {\bf Descripción} & El sistema deberá permanecer operativo durante todo el horario de trabajo de los empleados de lex.gal.  \\ \hline
\centering {\bf Importancia} & Vital \\ \hline
\centering {\bf Urgencia} & Inmediatamente \\ \hline
\end{tabular}
\caption{Requisito no funcional 04. Permanecer operativo durante el horario de trabajo.}
\label{enlaceNFR4}
\end{center}
\end{table}

\begin{table}[H]
\begin{center}
\begin{tabular}{|p{3cm}|p{10cm}|} \hline
\centering {\bf NFR-05} & Tratamiento de documentos en lenguajes de etiquetas  \\ \hline\hline
\centering {\bf Versión} & 1.0 (05/06/2022) \\ \hline
\centering {\bf Descripción} & El sistema deberá restringir el tratamiento de documentos a documentos escritos en lenguajes de etiquetas, como HTML o XML. \\ \hline
\centering {\bf Importancia} & Vital \\ \hline
\centering {\bf Urgencia} & Inmediatamente \\ \hline
\end{tabular}
\caption{Requisito no funcional 05. Tratamiento de documentos en lenguajes de etiquetas.}
\label{enlaceNFR5}
\end{center}
\end{table}

\begin{table}[H]
\begin{center}
\begin{tabular}{|p{3cm}|p{10cm}|} \hline
\centering {\bf NFR-06} & Tiempo de aprendizaje por el usuario $<$ 1h  \\ \hline\hline
\centering {\bf Versión} & 1.0 (05/06/2022) \\ \hline
\centering {\bf Descripción} & El sistema deberá tener un proceso de aprendizaje rápido y sencillo para el usuario, que en todo caso deberá ser menor de 1 hora. \\ \hline
\centering {\bf Importancia} & Vital \\ \hline
\centering {\bf Urgencia} & Inmediatamente \\ \hline
\end{tabular}
\caption{Requisito no funcional 06. Tiempo de aprendizaje por el usuario $<$ 1h.}
\label{enlaceNFR6}
\end{center}
\end{table}

\begin{table}[H]
\begin{center}
\begin{tabular}{|p{3cm}|p{10cm}|} \hline
\centering {\bf NFR-07} & Proporcionar mensajes de apoyo al usuario  \\ \hline\hline
\centering {\bf Versión} & 1.0 (05/06/2022) \\ \hline
\centering {\bf Descripción} & El sistema deberá proporcionar mensajes de apoyo al usuario en todas las operaciones, con el objetivo de facilitar la navegación en la aplicación. \\ \hline
\centering {\bf Importancia} & Quedaría bien \\ \hline
\centering {\bf Urgencia} & Puede esperar \\ \hline
\end{tabular}
\caption{Requisito no funcional 07. Proporcionar mensajes de apoyo al usuario.}
\label{enlaceNFR7}
\end{center}
\end{table}

\begin{table}[H]
\begin{center}
\begin{tabular}{|p{3cm}|p{10cm}|} \hline
\centering {\bf NFR-08} & Uso de tecnologías web y servicios RESTful  \\ \hline\hline
\centering {\bf Versión} & 1.0 (05/06/2022) \\ \hline
\centering {\bf Descripción} & El sistema deberá emplear tecnologías web para realizar la aplicación, incluyendo HTML5, Tailwind como framework de CSS y React.js como framework de JavaScript. Además, todos los servicios serán RESTful, y serán implementados con Spring Boot. \\ \hline
\centering {\bf Importancia} & Vital \\ \hline
\centering {\bf Urgencia} & Inmediatamente \\ \hline
\end{tabular}
\caption{Requisito no funcional 08. Uso de tecnologías web y servicios RESTful.}
\label{enlaceNFR8}
\end{center}
\end{table}

\begin{table}[H]
\begin{center}
\begin{tabular}{|p{3cm}|p{10cm}|} \hline
\centering {\bf NFR-09} & Base de datos NoSQL  \\ \hline\hline
\centering {\bf Versión} & 1.0 (05/06/2022) \\ \hline
\centering {\bf Descripción} & El sistema deberá almacenar todos los datos en una base de datos NoSQL de MongoDB. \\ \hline
\centering {\bf Importancia} & Vital \\ \hline
\centering {\bf Urgencia} & Inmediatamente \\ \hline
\end{tabular}
\caption{Requisito no funcional 09. Base de datos NoSQL.}
\label{enlaceNFR9}
\end{center}
\end{table}

\begin{table}[H]
\begin{center}
\begin{tabular}{|p{3cm}|p{10cm}|} \hline
\centering {\bf NFR-10} & Generación de tokens JWT  \\ \hline\hline
\centering {\bf Versión} & 1.0 (05/06/2022) \\ \hline
\centering {\bf Descripción} & El sistema deberá generar tokens JWT cada vez que un usuario inicie sesión en la aplicación. De esta forma, los datos del usuario no quedan tan accesibles en la web, y será necesario utilizar un algoritmo  {\it hash} para desencriptarlos. En todo caso la contraseña no será almacenada en el token JWT. \\ \hline
\centering {\bf Importancia} & Vital \\ \hline
\centering {\bf Urgencia} & Inmediatamente \\ \hline
\end{tabular}
\caption{Requisito no funcional 10. Generación de tokens JWT.}
\label{enlaceNFR10}
\end{center}
\end{table}
\section{Requisitos de información}

Por último, los requisitos de información \cite{requisitos}, son aquellos requisitos que indican algún tipo de información, la cual es necesaria para el correcto funcionamiento del sistema.
\\

En el caso de esta aplicación, se han localizado dos: Usuario y Documento.
\\

Los requisitos de información pueden ser consultados en el \hyperref[APRequisitosInformacion]{Apéndice C.3. Requisitos de información}.
\section{Casos de uso}

Un caso de uso \cite{casodeuso} es cualquier acción o actividad que se realiza en el sistema. Representan diferentes operaciones que se realizan en la aplicación.
\\

Los casos de uso pueden ser consultados en el \hyperref[APCasosUso]{Apéndice C.4. Casos de uso}.
\cleardoublepage
\chapter{Especificación de requisitos}
\label{enlaceespecificacion}

Para proceder a especificar los requisitos se ha realizado un profundo análisis de los objetivos de la aplicación. Estos requisitos se han buscado teniendo en cuenta aspectos funcionales, de rendimiento, de seguridad, de información, etc.
\\

Las plantillas empleadas para las tablas de los requisitos y casos de uso se encuentran en los apéndices: la \hyperref[enlaceFRQX]{E.1. Plantilla de requisitos funcionales}, la \hyperref[enlaceNFRX]{E.2. Plantilla de requisitos no funcionales}, la \hyperref[enlaceIRQX]{E.3. Plantilla de requisitos de información} y la \hyperref[enlaceUCX]{E.4. Plantilla de casos de uso}.
Estas plantillas están inspiradas en las tablas empleadas en el REM \cite{rem} para la gestión de requisitos.
\\

La importancia de los requisitos será medida según el grado de la necesidad del requisito en la aplicación, siendo en el caso de este proyecto ``Vital" para los requisitos de mayor importancia, y, ``Quedaría bien" para los menos importantes. Por su parte, la urgencia, indicará aquellas partes que se han de implementar con mayor rapidez, siendo ``Inmediatamente" las que tienen una mayor prioridad, y, ``Puede esperar", para los requisitos que no son estrictamente necesarios.

\section{Requisitos funcionales}

Los requisitos funcionales \cite{requisitos} son aquellos que definen una funcionalidad del sistema de software o sus componentes, tanto para las entradas de datos como para el comportamiento en situaciones específicas.
\\

En el caso de esta aplicación se han localizado dos grupos de requisitos funcionales: los asociados con operaciones de usuarios, y los que se relacionan con operaciones asociadas a leyes.
\\

Los requisitos funcionales pueden ser consultados en el \hyperref[APRequisitosFuncionales]{Apéndice C.1. Requisitos funcionales}.
\section{Requisitos no funcionales}
\label{APRequisitosNoFuncionales}

\begin{table}[H]
\begin{center}
\begin{tabular}{|p{3cm}|p{10cm}|} \hline
\centering {\bf NFR-01} & Tiempo de respuesta $<$ 2 segundos  \\ \hline\hline
\centering {\bf Versión} & 1.0 (05/06/2022) \\ \hline
\centering {\bf Descripción} & El sistema deberá responder en menos de 2 segundos, sin tener en cuenta el tiempo que tarde la red en enviar la información de vuelta al usuario. \\ \hline
\centering {\bf Importancia} & Vital \\ \hline
\centering {\bf Urgencia} & Inmediatamente \\ \hline
\end{tabular}
\caption{Requisito no funcional 01. Tiempo de respuesta $<$ 2 segundos.}
\label{enlaceNFR1}
\end{center}
\end{table}

\begin{table}[H]
\begin{center}
\begin{tabular}{|p{3cm}|p{10cm}|} \hline
\centering {\bf NFR-02} & Los datos sensibles del usuario se almacenan encriptados  \\ \hline\hline
\centering {\bf Versión} & 1.0 (05/06/2022) \\ \hline
\centering {\bf Descripción} & El sistema deberá encriptar los datos sensibles del usuario como las contraseñas en la base de datos de lex.gal. Se utilizará un algoritmo SHA-256 para encriptar. \\ \hline
\centering {\bf Importancia} & Vital \\ \hline
\centering {\bf Urgencia} & Inmediatamente \\ \hline
\end{tabular}
\caption{Requisito no funcional 02. Los datos sensibles del usuario se almacenan encriptados.}
\label{enlaceNFR2}
\end{center}
\end{table}

\begin{table}[H]
\begin{center}
\begin{tabular}{|p{3cm}|p{10cm}|} \hline
\centering {\bf NFR-03} & No permitir que un rol realice operaciones de otro  \\ \hline\hline
\centering {\bf Versión} & 1.0 (05/06/2022) \\ \hline
\centering {\bf Descripción} & El sistema deberá restringir aquellas operaciones exclusivas del administrador para el resto de usuarios. \\ \hline
\centering {\bf Importancia} & Vital \\ \hline
\centering {\bf Urgencia} & Inmediatamente \\ \hline
\end{tabular}
\caption{Requisito no funcional 03. No permitir que un rol realice operaciones de otro.}
\label{enlaceNFR3}
\end{center}
\end{table}

\begin{table}[H]
\begin{center}
\begin{tabular}{|p{3cm}|p{10cm}|} \hline
\centering {\bf NFR-04} & Permanecer operativo durante el horario de trabajo \\ \hline\hline
\centering {\bf Versión} & 1.0 (05/06/2022) \\ \hline
\centering {\bf Descripción} & El sistema deberá permanecer operativo durante todo el horario de trabajo de los empleados de lex.gal.  \\ \hline
\centering {\bf Importancia} & Vital \\ \hline
\centering {\bf Urgencia} & Inmediatamente \\ \hline
\end{tabular}
\caption{Requisito no funcional 04. Permanecer operativo durante el horario de trabajo.}
\label{enlaceNFR4}
\end{center}
\end{table}

\begin{table}[H]
\begin{center}
\begin{tabular}{|p{3cm}|p{10cm}|} \hline
\centering {\bf NFR-05} & Tratamiento de documentos en lenguajes de etiquetas  \\ \hline\hline
\centering {\bf Versión} & 1.0 (05/06/2022) \\ \hline
\centering {\bf Descripción} & El sistema deberá restringir el tratamiento de documentos a documentos escritos en lenguajes de etiquetas, como HTML o XML. \\ \hline
\centering {\bf Importancia} & Vital \\ \hline
\centering {\bf Urgencia} & Inmediatamente \\ \hline
\end{tabular}
\caption{Requisito no funcional 05. Tratamiento de documentos en lenguajes de etiquetas.}
\label{enlaceNFR5}
\end{center}
\end{table}

\begin{table}[H]
\begin{center}
\begin{tabular}{|p{3cm}|p{10cm}|} \hline
\centering {\bf NFR-06} & Tiempo de aprendizaje por el usuario $<$ 1h  \\ \hline\hline
\centering {\bf Versión} & 1.0 (05/06/2022) \\ \hline
\centering {\bf Descripción} & El sistema deberá tener un proceso de aprendizaje rápido y sencillo para el usuario, que en todo caso deberá ser menor de 1 hora. \\ \hline
\centering {\bf Importancia} & Vital \\ \hline
\centering {\bf Urgencia} & Inmediatamente \\ \hline
\end{tabular}
\caption{Requisito no funcional 06. Tiempo de aprendizaje por el usuario $<$ 1h.}
\label{enlaceNFR6}
\end{center}
\end{table}

\begin{table}[H]
\begin{center}
\begin{tabular}{|p{3cm}|p{10cm}|} \hline
\centering {\bf NFR-07} & Proporcionar mensajes de apoyo al usuario  \\ \hline\hline
\centering {\bf Versión} & 1.0 (05/06/2022) \\ \hline
\centering {\bf Descripción} & El sistema deberá proporcionar mensajes de apoyo al usuario en todas las operaciones, con el objetivo de facilitar la navegación en la aplicación. \\ \hline
\centering {\bf Importancia} & Quedaría bien \\ \hline
\centering {\bf Urgencia} & Puede esperar \\ \hline
\end{tabular}
\caption{Requisito no funcional 07. Proporcionar mensajes de apoyo al usuario.}
\label{enlaceNFR7}
\end{center}
\end{table}

\begin{table}[H]
\begin{center}
\begin{tabular}{|p{3cm}|p{10cm}|} \hline
\centering {\bf NFR-08} & Uso de tecnologías web y servicios RESTful  \\ \hline\hline
\centering {\bf Versión} & 1.0 (05/06/2022) \\ \hline
\centering {\bf Descripción} & El sistema deberá emplear tecnologías web para realizar la aplicación, incluyendo HTML5, Tailwind como framework de CSS y React.js como framework de JavaScript. Además, todos los servicios serán RESTful, y serán implementados con Spring Boot. \\ \hline
\centering {\bf Importancia} & Vital \\ \hline
\centering {\bf Urgencia} & Inmediatamente \\ \hline
\end{tabular}
\caption{Requisito no funcional 08. Uso de tecnologías web y servicios RESTful.}
\label{enlaceNFR8}
\end{center}
\end{table}

\begin{table}[H]
\begin{center}
\begin{tabular}{|p{3cm}|p{10cm}|} \hline
\centering {\bf NFR-09} & Base de datos NoSQL  \\ \hline\hline
\centering {\bf Versión} & 1.0 (05/06/2022) \\ \hline
\centering {\bf Descripción} & El sistema deberá almacenar todos los datos en una base de datos NoSQL de MongoDB. \\ \hline
\centering {\bf Importancia} & Vital \\ \hline
\centering {\bf Urgencia} & Inmediatamente \\ \hline
\end{tabular}
\caption{Requisito no funcional 09. Base de datos NoSQL.}
\label{enlaceNFR9}
\end{center}
\end{table}

\begin{table}[H]
\begin{center}
\begin{tabular}{|p{3cm}|p{10cm}|} \hline
\centering {\bf NFR-10} & Generación de tokens JWT  \\ \hline\hline
\centering {\bf Versión} & 1.0 (05/06/2022) \\ \hline
\centering {\bf Descripción} & El sistema deberá generar tokens JWT cada vez que un usuario inicie sesión en la aplicación. De esta forma, los datos del usuario no quedan tan accesibles en la web, y será necesario utilizar un algoritmo  {\it hash} para desencriptarlos. En todo caso la contraseña no será almacenada en el token JWT. \\ \hline
\centering {\bf Importancia} & Vital \\ \hline
\centering {\bf Urgencia} & Inmediatamente \\ \hline
\end{tabular}
\caption{Requisito no funcional 10. Generación de tokens JWT.}
\label{enlaceNFR10}
\end{center}
\end{table}
\section{Requisitos de información}

Por último, los requisitos de información \cite{requisitos}, son aquellos requisitos que indican algún tipo de información, la cual es necesaria para el correcto funcionamiento del sistema.
\\

En el caso de esta aplicación, se han localizado dos: Usuario y Documento.
\\

Los requisitos de información pueden ser consultados en el \hyperref[APRequisitosInformacion]{Apéndice C.3. Requisitos de información}.
\section{Casos de uso}

Un caso de uso \cite{casodeuso} es cualquier acción o actividad que se realiza en el sistema. Representan diferentes operaciones que se realizan en la aplicación.
\\

Los casos de uso pueden ser consultados en el \hyperref[APCasosUso]{Apéndice C.4. Casos de uso}.
\cleardoublepage
\chapter{Especificación de requisitos}
\label{enlaceespecificacion}

Para proceder a especificar los requisitos se ha realizado un profundo análisis de los objetivos de la aplicación. Estos requisitos se han buscado teniendo en cuenta aspectos funcionales, de rendimiento, de seguridad, de información, etc.
\\

Las plantillas empleadas para las tablas de los requisitos y casos de uso se encuentran en los apéndices: la \hyperref[enlaceFRQX]{E.1. Plantilla de requisitos funcionales}, la \hyperref[enlaceNFRX]{E.2. Plantilla de requisitos no funcionales}, la \hyperref[enlaceIRQX]{E.3. Plantilla de requisitos de información} y la \hyperref[enlaceUCX]{E.4. Plantilla de casos de uso}.
Estas plantillas están inspiradas en las tablas empleadas en el REM \cite{rem} para la gestión de requisitos.
\\

La importancia de los requisitos será medida según el grado de la necesidad del requisito en la aplicación, siendo en el caso de este proyecto ``Vital" para los requisitos de mayor importancia, y, ``Quedaría bien" para los menos importantes. Por su parte, la urgencia, indicará aquellas partes que se han de implementar con mayor rapidez, siendo ``Inmediatamente" las que tienen una mayor prioridad, y, ``Puede esperar", para los requisitos que no son estrictamente necesarios.

\section{Requisitos funcionales}

Los requisitos funcionales \cite{requisitos} son aquellos que definen una funcionalidad del sistema de software o sus componentes, tanto para las entradas de datos como para el comportamiento en situaciones específicas.
\\

En el caso de esta aplicación se han localizado dos grupos de requisitos funcionales: los asociados con operaciones de usuarios, y los que se relacionan con operaciones asociadas a leyes.
\\

Los requisitos funcionales pueden ser consultados en el \hyperref[APRequisitosFuncionales]{Apéndice C.1. Requisitos funcionales}.
\section{Requisitos no funcionales}
\label{APRequisitosNoFuncionales}

\begin{table}[H]
\begin{center}
\begin{tabular}{|p{3cm}|p{10cm}|} \hline
\centering {\bf NFR-01} & Tiempo de respuesta $<$ 2 segundos  \\ \hline\hline
\centering {\bf Versión} & 1.0 (05/06/2022) \\ \hline
\centering {\bf Descripción} & El sistema deberá responder en menos de 2 segundos, sin tener en cuenta el tiempo que tarde la red en enviar la información de vuelta al usuario. \\ \hline
\centering {\bf Importancia} & Vital \\ \hline
\centering {\bf Urgencia} & Inmediatamente \\ \hline
\end{tabular}
\caption{Requisito no funcional 01. Tiempo de respuesta $<$ 2 segundos.}
\label{enlaceNFR1}
\end{center}
\end{table}

\begin{table}[H]
\begin{center}
\begin{tabular}{|p{3cm}|p{10cm}|} \hline
\centering {\bf NFR-02} & Los datos sensibles del usuario se almacenan encriptados  \\ \hline\hline
\centering {\bf Versión} & 1.0 (05/06/2022) \\ \hline
\centering {\bf Descripción} & El sistema deberá encriptar los datos sensibles del usuario como las contraseñas en la base de datos de lex.gal. Se utilizará un algoritmo SHA-256 para encriptar. \\ \hline
\centering {\bf Importancia} & Vital \\ \hline
\centering {\bf Urgencia} & Inmediatamente \\ \hline
\end{tabular}
\caption{Requisito no funcional 02. Los datos sensibles del usuario se almacenan encriptados.}
\label{enlaceNFR2}
\end{center}
\end{table}

\begin{table}[H]
\begin{center}
\begin{tabular}{|p{3cm}|p{10cm}|} \hline
\centering {\bf NFR-03} & No permitir que un rol realice operaciones de otro  \\ \hline\hline
\centering {\bf Versión} & 1.0 (05/06/2022) \\ \hline
\centering {\bf Descripción} & El sistema deberá restringir aquellas operaciones exclusivas del administrador para el resto de usuarios. \\ \hline
\centering {\bf Importancia} & Vital \\ \hline
\centering {\bf Urgencia} & Inmediatamente \\ \hline
\end{tabular}
\caption{Requisito no funcional 03. No permitir que un rol realice operaciones de otro.}
\label{enlaceNFR3}
\end{center}
\end{table}

\begin{table}[H]
\begin{center}
\begin{tabular}{|p{3cm}|p{10cm}|} \hline
\centering {\bf NFR-04} & Permanecer operativo durante el horario de trabajo \\ \hline\hline
\centering {\bf Versión} & 1.0 (05/06/2022) \\ \hline
\centering {\bf Descripción} & El sistema deberá permanecer operativo durante todo el horario de trabajo de los empleados de lex.gal.  \\ \hline
\centering {\bf Importancia} & Vital \\ \hline
\centering {\bf Urgencia} & Inmediatamente \\ \hline
\end{tabular}
\caption{Requisito no funcional 04. Permanecer operativo durante el horario de trabajo.}
\label{enlaceNFR4}
\end{center}
\end{table}

\begin{table}[H]
\begin{center}
\begin{tabular}{|p{3cm}|p{10cm}|} \hline
\centering {\bf NFR-05} & Tratamiento de documentos en lenguajes de etiquetas  \\ \hline\hline
\centering {\bf Versión} & 1.0 (05/06/2022) \\ \hline
\centering {\bf Descripción} & El sistema deberá restringir el tratamiento de documentos a documentos escritos en lenguajes de etiquetas, como HTML o XML. \\ \hline
\centering {\bf Importancia} & Vital \\ \hline
\centering {\bf Urgencia} & Inmediatamente \\ \hline
\end{tabular}
\caption{Requisito no funcional 05. Tratamiento de documentos en lenguajes de etiquetas.}
\label{enlaceNFR5}
\end{center}
\end{table}

\begin{table}[H]
\begin{center}
\begin{tabular}{|p{3cm}|p{10cm}|} \hline
\centering {\bf NFR-06} & Tiempo de aprendizaje por el usuario $<$ 1h  \\ \hline\hline
\centering {\bf Versión} & 1.0 (05/06/2022) \\ \hline
\centering {\bf Descripción} & El sistema deberá tener un proceso de aprendizaje rápido y sencillo para el usuario, que en todo caso deberá ser menor de 1 hora. \\ \hline
\centering {\bf Importancia} & Vital \\ \hline
\centering {\bf Urgencia} & Inmediatamente \\ \hline
\end{tabular}
\caption{Requisito no funcional 06. Tiempo de aprendizaje por el usuario $<$ 1h.}
\label{enlaceNFR6}
\end{center}
\end{table}

\begin{table}[H]
\begin{center}
\begin{tabular}{|p{3cm}|p{10cm}|} \hline
\centering {\bf NFR-07} & Proporcionar mensajes de apoyo al usuario  \\ \hline\hline
\centering {\bf Versión} & 1.0 (05/06/2022) \\ \hline
\centering {\bf Descripción} & El sistema deberá proporcionar mensajes de apoyo al usuario en todas las operaciones, con el objetivo de facilitar la navegación en la aplicación. \\ \hline
\centering {\bf Importancia} & Quedaría bien \\ \hline
\centering {\bf Urgencia} & Puede esperar \\ \hline
\end{tabular}
\caption{Requisito no funcional 07. Proporcionar mensajes de apoyo al usuario.}
\label{enlaceNFR7}
\end{center}
\end{table}

\begin{table}[H]
\begin{center}
\begin{tabular}{|p{3cm}|p{10cm}|} \hline
\centering {\bf NFR-08} & Uso de tecnologías web y servicios RESTful  \\ \hline\hline
\centering {\bf Versión} & 1.0 (05/06/2022) \\ \hline
\centering {\bf Descripción} & El sistema deberá emplear tecnologías web para realizar la aplicación, incluyendo HTML5, Tailwind como framework de CSS y React.js como framework de JavaScript. Además, todos los servicios serán RESTful, y serán implementados con Spring Boot. \\ \hline
\centering {\bf Importancia} & Vital \\ \hline
\centering {\bf Urgencia} & Inmediatamente \\ \hline
\end{tabular}
\caption{Requisito no funcional 08. Uso de tecnologías web y servicios RESTful.}
\label{enlaceNFR8}
\end{center}
\end{table}

\begin{table}[H]
\begin{center}
\begin{tabular}{|p{3cm}|p{10cm}|} \hline
\centering {\bf NFR-09} & Base de datos NoSQL  \\ \hline\hline
\centering {\bf Versión} & 1.0 (05/06/2022) \\ \hline
\centering {\bf Descripción} & El sistema deberá almacenar todos los datos en una base de datos NoSQL de MongoDB. \\ \hline
\centering {\bf Importancia} & Vital \\ \hline
\centering {\bf Urgencia} & Inmediatamente \\ \hline
\end{tabular}
\caption{Requisito no funcional 09. Base de datos NoSQL.}
\label{enlaceNFR9}
\end{center}
\end{table}

\begin{table}[H]
\begin{center}
\begin{tabular}{|p{3cm}|p{10cm}|} \hline
\centering {\bf NFR-10} & Generación de tokens JWT  \\ \hline\hline
\centering {\bf Versión} & 1.0 (05/06/2022) \\ \hline
\centering {\bf Descripción} & El sistema deberá generar tokens JWT cada vez que un usuario inicie sesión en la aplicación. De esta forma, los datos del usuario no quedan tan accesibles en la web, y será necesario utilizar un algoritmo  {\it hash} para desencriptarlos. En todo caso la contraseña no será almacenada en el token JWT. \\ \hline
\centering {\bf Importancia} & Vital \\ \hline
\centering {\bf Urgencia} & Inmediatamente \\ \hline
\end{tabular}
\caption{Requisito no funcional 10. Generación de tokens JWT.}
\label{enlaceNFR10}
\end{center}
\end{table}
\section{Requisitos de información}

Por último, los requisitos de información \cite{requisitos}, son aquellos requisitos que indican algún tipo de información, la cual es necesaria para el correcto funcionamiento del sistema.
\\

En el caso de esta aplicación, se han localizado dos: Usuario y Documento.
\\

Los requisitos de información pueden ser consultados en el \hyperref[APRequisitosInformacion]{Apéndice C.3. Requisitos de información}.
\section{Casos de uso}

Un caso de uso \cite{casodeuso} es cualquier acción o actividad que se realiza en el sistema. Representan diferentes operaciones que se realizan en la aplicación.
\\

Los casos de uso pueden ser consultados en el \hyperref[APCasosUso]{Apéndice C.4. Casos de uso}.
\cleardoublepage
\chapter{Especificación de requisitos}
\label{enlaceespecificacion}

Para proceder a especificar los requisitos se ha realizado un profundo análisis de los objetivos de la aplicación. Estos requisitos se han buscado teniendo en cuenta aspectos funcionales, de rendimiento, de seguridad, de información, etc.
\\

Las plantillas empleadas para las tablas de los requisitos y casos de uso se encuentran en los apéndices: la \hyperref[enlaceFRQX]{E.1. Plantilla de requisitos funcionales}, la \hyperref[enlaceNFRX]{E.2. Plantilla de requisitos no funcionales}, la \hyperref[enlaceIRQX]{E.3. Plantilla de requisitos de información} y la \hyperref[enlaceUCX]{E.4. Plantilla de casos de uso}.
Estas plantillas están inspiradas en las tablas empleadas en el REM \cite{rem} para la gestión de requisitos.
\\

La importancia de los requisitos será medida según el grado de la necesidad del requisito en la aplicación, siendo en el caso de este proyecto ``Vital" para los requisitos de mayor importancia, y, ``Quedaría bien" para los menos importantes. Por su parte, la urgencia, indicará aquellas partes que se han de implementar con mayor rapidez, siendo ``Inmediatamente" las que tienen una mayor prioridad, y, ``Puede esperar", para los requisitos que no son estrictamente necesarios.

\section{Requisitos funcionales}

Los requisitos funcionales \cite{requisitos} son aquellos que definen una funcionalidad del sistema de software o sus componentes, tanto para las entradas de datos como para el comportamiento en situaciones específicas.
\\

En el caso de esta aplicación se han localizado dos grupos de requisitos funcionales: los asociados con operaciones de usuarios, y los que se relacionan con operaciones asociadas a leyes.
\\

Los requisitos funcionales pueden ser consultados en el \hyperref[APRequisitosFuncionales]{Apéndice C.1. Requisitos funcionales}.
\section{Requisitos no funcionales}
\label{APRequisitosNoFuncionales}

\begin{table}[H]
\begin{center}
\begin{tabular}{|p{3cm}|p{10cm}|} \hline
\centering {\bf NFR-01} & Tiempo de respuesta $<$ 2 segundos  \\ \hline\hline
\centering {\bf Versión} & 1.0 (05/06/2022) \\ \hline
\centering {\bf Descripción} & El sistema deberá responder en menos de 2 segundos, sin tener en cuenta el tiempo que tarde la red en enviar la información de vuelta al usuario. \\ \hline
\centering {\bf Importancia} & Vital \\ \hline
\centering {\bf Urgencia} & Inmediatamente \\ \hline
\end{tabular}
\caption{Requisito no funcional 01. Tiempo de respuesta $<$ 2 segundos.}
\label{enlaceNFR1}
\end{center}
\end{table}

\begin{table}[H]
\begin{center}
\begin{tabular}{|p{3cm}|p{10cm}|} \hline
\centering {\bf NFR-02} & Los datos sensibles del usuario se almacenan encriptados  \\ \hline\hline
\centering {\bf Versión} & 1.0 (05/06/2022) \\ \hline
\centering {\bf Descripción} & El sistema deberá encriptar los datos sensibles del usuario como las contraseñas en la base de datos de lex.gal. Se utilizará un algoritmo SHA-256 para encriptar. \\ \hline
\centering {\bf Importancia} & Vital \\ \hline
\centering {\bf Urgencia} & Inmediatamente \\ \hline
\end{tabular}
\caption{Requisito no funcional 02. Los datos sensibles del usuario se almacenan encriptados.}
\label{enlaceNFR2}
\end{center}
\end{table}

\begin{table}[H]
\begin{center}
\begin{tabular}{|p{3cm}|p{10cm}|} \hline
\centering {\bf NFR-03} & No permitir que un rol realice operaciones de otro  \\ \hline\hline
\centering {\bf Versión} & 1.0 (05/06/2022) \\ \hline
\centering {\bf Descripción} & El sistema deberá restringir aquellas operaciones exclusivas del administrador para el resto de usuarios. \\ \hline
\centering {\bf Importancia} & Vital \\ \hline
\centering {\bf Urgencia} & Inmediatamente \\ \hline
\end{tabular}
\caption{Requisito no funcional 03. No permitir que un rol realice operaciones de otro.}
\label{enlaceNFR3}
\end{center}
\end{table}

\begin{table}[H]
\begin{center}
\begin{tabular}{|p{3cm}|p{10cm}|} \hline
\centering {\bf NFR-04} & Permanecer operativo durante el horario de trabajo \\ \hline\hline
\centering {\bf Versión} & 1.0 (05/06/2022) \\ \hline
\centering {\bf Descripción} & El sistema deberá permanecer operativo durante todo el horario de trabajo de los empleados de lex.gal.  \\ \hline
\centering {\bf Importancia} & Vital \\ \hline
\centering {\bf Urgencia} & Inmediatamente \\ \hline
\end{tabular}
\caption{Requisito no funcional 04. Permanecer operativo durante el horario de trabajo.}
\label{enlaceNFR4}
\end{center}
\end{table}

\begin{table}[H]
\begin{center}
\begin{tabular}{|p{3cm}|p{10cm}|} \hline
\centering {\bf NFR-05} & Tratamiento de documentos en lenguajes de etiquetas  \\ \hline\hline
\centering {\bf Versión} & 1.0 (05/06/2022) \\ \hline
\centering {\bf Descripción} & El sistema deberá restringir el tratamiento de documentos a documentos escritos en lenguajes de etiquetas, como HTML o XML. \\ \hline
\centering {\bf Importancia} & Vital \\ \hline
\centering {\bf Urgencia} & Inmediatamente \\ \hline
\end{tabular}
\caption{Requisito no funcional 05. Tratamiento de documentos en lenguajes de etiquetas.}
\label{enlaceNFR5}
\end{center}
\end{table}

\begin{table}[H]
\begin{center}
\begin{tabular}{|p{3cm}|p{10cm}|} \hline
\centering {\bf NFR-06} & Tiempo de aprendizaje por el usuario $<$ 1h  \\ \hline\hline
\centering {\bf Versión} & 1.0 (05/06/2022) \\ \hline
\centering {\bf Descripción} & El sistema deberá tener un proceso de aprendizaje rápido y sencillo para el usuario, que en todo caso deberá ser menor de 1 hora. \\ \hline
\centering {\bf Importancia} & Vital \\ \hline
\centering {\bf Urgencia} & Inmediatamente \\ \hline
\end{tabular}
\caption{Requisito no funcional 06. Tiempo de aprendizaje por el usuario $<$ 1h.}
\label{enlaceNFR6}
\end{center}
\end{table}

\begin{table}[H]
\begin{center}
\begin{tabular}{|p{3cm}|p{10cm}|} \hline
\centering {\bf NFR-07} & Proporcionar mensajes de apoyo al usuario  \\ \hline\hline
\centering {\bf Versión} & 1.0 (05/06/2022) \\ \hline
\centering {\bf Descripción} & El sistema deberá proporcionar mensajes de apoyo al usuario en todas las operaciones, con el objetivo de facilitar la navegación en la aplicación. \\ \hline
\centering {\bf Importancia} & Quedaría bien \\ \hline
\centering {\bf Urgencia} & Puede esperar \\ \hline
\end{tabular}
\caption{Requisito no funcional 07. Proporcionar mensajes de apoyo al usuario.}
\label{enlaceNFR7}
\end{center}
\end{table}

\begin{table}[H]
\begin{center}
\begin{tabular}{|p{3cm}|p{10cm}|} \hline
\centering {\bf NFR-08} & Uso de tecnologías web y servicios RESTful  \\ \hline\hline
\centering {\bf Versión} & 1.0 (05/06/2022) \\ \hline
\centering {\bf Descripción} & El sistema deberá emplear tecnologías web para realizar la aplicación, incluyendo HTML5, Tailwind como framework de CSS y React.js como framework de JavaScript. Además, todos los servicios serán RESTful, y serán implementados con Spring Boot. \\ \hline
\centering {\bf Importancia} & Vital \\ \hline
\centering {\bf Urgencia} & Inmediatamente \\ \hline
\end{tabular}
\caption{Requisito no funcional 08. Uso de tecnologías web y servicios RESTful.}
\label{enlaceNFR8}
\end{center}
\end{table}

\begin{table}[H]
\begin{center}
\begin{tabular}{|p{3cm}|p{10cm}|} \hline
\centering {\bf NFR-09} & Base de datos NoSQL  \\ \hline\hline
\centering {\bf Versión} & 1.0 (05/06/2022) \\ \hline
\centering {\bf Descripción} & El sistema deberá almacenar todos los datos en una base de datos NoSQL de MongoDB. \\ \hline
\centering {\bf Importancia} & Vital \\ \hline
\centering {\bf Urgencia} & Inmediatamente \\ \hline
\end{tabular}
\caption{Requisito no funcional 09. Base de datos NoSQL.}
\label{enlaceNFR9}
\end{center}
\end{table}

\begin{table}[H]
\begin{center}
\begin{tabular}{|p{3cm}|p{10cm}|} \hline
\centering {\bf NFR-10} & Generación de tokens JWT  \\ \hline\hline
\centering {\bf Versión} & 1.0 (05/06/2022) \\ \hline
\centering {\bf Descripción} & El sistema deberá generar tokens JWT cada vez que un usuario inicie sesión en la aplicación. De esta forma, los datos del usuario no quedan tan accesibles en la web, y será necesario utilizar un algoritmo  {\it hash} para desencriptarlos. En todo caso la contraseña no será almacenada en el token JWT. \\ \hline
\centering {\bf Importancia} & Vital \\ \hline
\centering {\bf Urgencia} & Inmediatamente \\ \hline
\end{tabular}
\caption{Requisito no funcional 10. Generación de tokens JWT.}
\label{enlaceNFR10}
\end{center}
\end{table}
\section{Requisitos de información}

Por último, los requisitos de información \cite{requisitos}, son aquellos requisitos que indican algún tipo de información, la cual es necesaria para el correcto funcionamiento del sistema.
\\

En el caso de esta aplicación, se han localizado dos: Usuario y Documento.
\\

Los requisitos de información pueden ser consultados en el \hyperref[APRequisitosInformacion]{Apéndice C.3. Requisitos de información}.
\section{Casos de uso}

Un caso de uso \cite{casodeuso} es cualquier acción o actividad que se realiza en el sistema. Representan diferentes operaciones que se realizan en la aplicación.
\\

Los casos de uso pueden ser consultados en el \hyperref[APCasosUso]{Apéndice C.4. Casos de uso}.
\cleardoublepage
\chapter{Especificación de requisitos}
\label{enlaceespecificacion}

Para proceder a especificar los requisitos se ha realizado un profundo análisis de los objetivos de la aplicación. Estos requisitos se han buscado teniendo en cuenta aspectos funcionales, de rendimiento, de seguridad, de información, etc.
\\

Las plantillas empleadas para las tablas de los requisitos y casos de uso se encuentran en los apéndices: la \hyperref[enlaceFRQX]{E.1. Plantilla de requisitos funcionales}, la \hyperref[enlaceNFRX]{E.2. Plantilla de requisitos no funcionales}, la \hyperref[enlaceIRQX]{E.3. Plantilla de requisitos de información} y la \hyperref[enlaceUCX]{E.4. Plantilla de casos de uso}.
Estas plantillas están inspiradas en las tablas empleadas en el REM \cite{rem} para la gestión de requisitos.
\\

La importancia de los requisitos será medida según el grado de la necesidad del requisito en la aplicación, siendo en el caso de este proyecto ``Vital" para los requisitos de mayor importancia, y, ``Quedaría bien" para los menos importantes. Por su parte, la urgencia, indicará aquellas partes que se han de implementar con mayor rapidez, siendo ``Inmediatamente" las que tienen una mayor prioridad, y, ``Puede esperar", para los requisitos que no son estrictamente necesarios.

\section{Requisitos funcionales}

Los requisitos funcionales \cite{requisitos} son aquellos que definen una funcionalidad del sistema de software o sus componentes, tanto para las entradas de datos como para el comportamiento en situaciones específicas.
\\

En el caso de esta aplicación se han localizado dos grupos de requisitos funcionales: los asociados con operaciones de usuarios, y los que se relacionan con operaciones asociadas a leyes.
\\

Los requisitos funcionales pueden ser consultados en el \hyperref[APRequisitosFuncionales]{Apéndice C.1. Requisitos funcionales}.
\section{Requisitos no funcionales}
\label{APRequisitosNoFuncionales}

\begin{table}[H]
\begin{center}
\begin{tabular}{|p{3cm}|p{10cm}|} \hline
\centering {\bf NFR-01} & Tiempo de respuesta $<$ 2 segundos  \\ \hline\hline
\centering {\bf Versión} & 1.0 (05/06/2022) \\ \hline
\centering {\bf Descripción} & El sistema deberá responder en menos de 2 segundos, sin tener en cuenta el tiempo que tarde la red en enviar la información de vuelta al usuario. \\ \hline
\centering {\bf Importancia} & Vital \\ \hline
\centering {\bf Urgencia} & Inmediatamente \\ \hline
\end{tabular}
\caption{Requisito no funcional 01. Tiempo de respuesta $<$ 2 segundos.}
\label{enlaceNFR1}
\end{center}
\end{table}

\begin{table}[H]
\begin{center}
\begin{tabular}{|p{3cm}|p{10cm}|} \hline
\centering {\bf NFR-02} & Los datos sensibles del usuario se almacenan encriptados  \\ \hline\hline
\centering {\bf Versión} & 1.0 (05/06/2022) \\ \hline
\centering {\bf Descripción} & El sistema deberá encriptar los datos sensibles del usuario como las contraseñas en la base de datos de lex.gal. Se utilizará un algoritmo SHA-256 para encriptar. \\ \hline
\centering {\bf Importancia} & Vital \\ \hline
\centering {\bf Urgencia} & Inmediatamente \\ \hline
\end{tabular}
\caption{Requisito no funcional 02. Los datos sensibles del usuario se almacenan encriptados.}
\label{enlaceNFR2}
\end{center}
\end{table}

\begin{table}[H]
\begin{center}
\begin{tabular}{|p{3cm}|p{10cm}|} \hline
\centering {\bf NFR-03} & No permitir que un rol realice operaciones de otro  \\ \hline\hline
\centering {\bf Versión} & 1.0 (05/06/2022) \\ \hline
\centering {\bf Descripción} & El sistema deberá restringir aquellas operaciones exclusivas del administrador para el resto de usuarios. \\ \hline
\centering {\bf Importancia} & Vital \\ \hline
\centering {\bf Urgencia} & Inmediatamente \\ \hline
\end{tabular}
\caption{Requisito no funcional 03. No permitir que un rol realice operaciones de otro.}
\label{enlaceNFR3}
\end{center}
\end{table}

\begin{table}[H]
\begin{center}
\begin{tabular}{|p{3cm}|p{10cm}|} \hline
\centering {\bf NFR-04} & Permanecer operativo durante el horario de trabajo \\ \hline\hline
\centering {\bf Versión} & 1.0 (05/06/2022) \\ \hline
\centering {\bf Descripción} & El sistema deberá permanecer operativo durante todo el horario de trabajo de los empleados de lex.gal.  \\ \hline
\centering {\bf Importancia} & Vital \\ \hline
\centering {\bf Urgencia} & Inmediatamente \\ \hline
\end{tabular}
\caption{Requisito no funcional 04. Permanecer operativo durante el horario de trabajo.}
\label{enlaceNFR4}
\end{center}
\end{table}

\begin{table}[H]
\begin{center}
\begin{tabular}{|p{3cm}|p{10cm}|} \hline
\centering {\bf NFR-05} & Tratamiento de documentos en lenguajes de etiquetas  \\ \hline\hline
\centering {\bf Versión} & 1.0 (05/06/2022) \\ \hline
\centering {\bf Descripción} & El sistema deberá restringir el tratamiento de documentos a documentos escritos en lenguajes de etiquetas, como HTML o XML. \\ \hline
\centering {\bf Importancia} & Vital \\ \hline
\centering {\bf Urgencia} & Inmediatamente \\ \hline
\end{tabular}
\caption{Requisito no funcional 05. Tratamiento de documentos en lenguajes de etiquetas.}
\label{enlaceNFR5}
\end{center}
\end{table}

\begin{table}[H]
\begin{center}
\begin{tabular}{|p{3cm}|p{10cm}|} \hline
\centering {\bf NFR-06} & Tiempo de aprendizaje por el usuario $<$ 1h  \\ \hline\hline
\centering {\bf Versión} & 1.0 (05/06/2022) \\ \hline
\centering {\bf Descripción} & El sistema deberá tener un proceso de aprendizaje rápido y sencillo para el usuario, que en todo caso deberá ser menor de 1 hora. \\ \hline
\centering {\bf Importancia} & Vital \\ \hline
\centering {\bf Urgencia} & Inmediatamente \\ \hline
\end{tabular}
\caption{Requisito no funcional 06. Tiempo de aprendizaje por el usuario $<$ 1h.}
\label{enlaceNFR6}
\end{center}
\end{table}

\begin{table}[H]
\begin{center}
\begin{tabular}{|p{3cm}|p{10cm}|} \hline
\centering {\bf NFR-07} & Proporcionar mensajes de apoyo al usuario  \\ \hline\hline
\centering {\bf Versión} & 1.0 (05/06/2022) \\ \hline
\centering {\bf Descripción} & El sistema deberá proporcionar mensajes de apoyo al usuario en todas las operaciones, con el objetivo de facilitar la navegación en la aplicación. \\ \hline
\centering {\bf Importancia} & Quedaría bien \\ \hline
\centering {\bf Urgencia} & Puede esperar \\ \hline
\end{tabular}
\caption{Requisito no funcional 07. Proporcionar mensajes de apoyo al usuario.}
\label{enlaceNFR7}
\end{center}
\end{table}

\begin{table}[H]
\begin{center}
\begin{tabular}{|p{3cm}|p{10cm}|} \hline
\centering {\bf NFR-08} & Uso de tecnologías web y servicios RESTful  \\ \hline\hline
\centering {\bf Versión} & 1.0 (05/06/2022) \\ \hline
\centering {\bf Descripción} & El sistema deberá emplear tecnologías web para realizar la aplicación, incluyendo HTML5, Tailwind como framework de CSS y React.js como framework de JavaScript. Además, todos los servicios serán RESTful, y serán implementados con Spring Boot. \\ \hline
\centering {\bf Importancia} & Vital \\ \hline
\centering {\bf Urgencia} & Inmediatamente \\ \hline
\end{tabular}
\caption{Requisito no funcional 08. Uso de tecnologías web y servicios RESTful.}
\label{enlaceNFR8}
\end{center}
\end{table}

\begin{table}[H]
\begin{center}
\begin{tabular}{|p{3cm}|p{10cm}|} \hline
\centering {\bf NFR-09} & Base de datos NoSQL  \\ \hline\hline
\centering {\bf Versión} & 1.0 (05/06/2022) \\ \hline
\centering {\bf Descripción} & El sistema deberá almacenar todos los datos en una base de datos NoSQL de MongoDB. \\ \hline
\centering {\bf Importancia} & Vital \\ \hline
\centering {\bf Urgencia} & Inmediatamente \\ \hline
\end{tabular}
\caption{Requisito no funcional 09. Base de datos NoSQL.}
\label{enlaceNFR9}
\end{center}
\end{table}

\begin{table}[H]
\begin{center}
\begin{tabular}{|p{3cm}|p{10cm}|} \hline
\centering {\bf NFR-10} & Generación de tokens JWT  \\ \hline\hline
\centering {\bf Versión} & 1.0 (05/06/2022) \\ \hline
\centering {\bf Descripción} & El sistema deberá generar tokens JWT cada vez que un usuario inicie sesión en la aplicación. De esta forma, los datos del usuario no quedan tan accesibles en la web, y será necesario utilizar un algoritmo  {\it hash} para desencriptarlos. En todo caso la contraseña no será almacenada en el token JWT. \\ \hline
\centering {\bf Importancia} & Vital \\ \hline
\centering {\bf Urgencia} & Inmediatamente \\ \hline
\end{tabular}
\caption{Requisito no funcional 10. Generación de tokens JWT.}
\label{enlaceNFR10}
\end{center}
\end{table}
\section{Requisitos de información}

Por último, los requisitos de información \cite{requisitos}, son aquellos requisitos que indican algún tipo de información, la cual es necesaria para el correcto funcionamiento del sistema.
\\

En el caso de esta aplicación, se han localizado dos: Usuario y Documento.
\\

Los requisitos de información pueden ser consultados en el \hyperref[APRequisitosInformacion]{Apéndice C.3. Requisitos de información}.
\section{Casos de uso}

Un caso de uso \cite{casodeuso} es cualquier acción o actividad que se realiza en el sistema. Representan diferentes operaciones que se realizan en la aplicación.
\\

Los casos de uso pueden ser consultados en el \hyperref[APCasosUso]{Apéndice C.4. Casos de uso}.
\cleardoublepage

% Aquí empezan os apéndices
\appendix
\cleardoublepage
\chapter{Especificación de requisitos}
\label{enlaceespecificacion}

Para proceder a especificar los requisitos se ha realizado un profundo análisis de los objetivos de la aplicación. Estos requisitos se han buscado teniendo en cuenta aspectos funcionales, de rendimiento, de seguridad, de información, etc.
\\

Las plantillas empleadas para las tablas de los requisitos y casos de uso se encuentran en los apéndices: la \hyperref[enlaceFRQX]{E.1. Plantilla de requisitos funcionales}, la \hyperref[enlaceNFRX]{E.2. Plantilla de requisitos no funcionales}, la \hyperref[enlaceIRQX]{E.3. Plantilla de requisitos de información} y la \hyperref[enlaceUCX]{E.4. Plantilla de casos de uso}.
Estas plantillas están inspiradas en las tablas empleadas en el REM \cite{rem} para la gestión de requisitos.
\\

La importancia de los requisitos será medida según el grado de la necesidad del requisito en la aplicación, siendo en el caso de este proyecto ``Vital" para los requisitos de mayor importancia, y, ``Quedaría bien" para los menos importantes. Por su parte, la urgencia, indicará aquellas partes que se han de implementar con mayor rapidez, siendo ``Inmediatamente" las que tienen una mayor prioridad, y, ``Puede esperar", para los requisitos que no son estrictamente necesarios.

\section{Requisitos funcionales}

Los requisitos funcionales \cite{requisitos} son aquellos que definen una funcionalidad del sistema de software o sus componentes, tanto para las entradas de datos como para el comportamiento en situaciones específicas.
\\

En el caso de esta aplicación se han localizado dos grupos de requisitos funcionales: los asociados con operaciones de usuarios, y los que se relacionan con operaciones asociadas a leyes.
\\

Los requisitos funcionales pueden ser consultados en el \hyperref[APRequisitosFuncionales]{Apéndice C.1. Requisitos funcionales}.
\section{Requisitos no funcionales}
\label{APRequisitosNoFuncionales}

\begin{table}[H]
\begin{center}
\begin{tabular}{|p{3cm}|p{10cm}|} \hline
\centering {\bf NFR-01} & Tiempo de respuesta $<$ 2 segundos  \\ \hline\hline
\centering {\bf Versión} & 1.0 (05/06/2022) \\ \hline
\centering {\bf Descripción} & El sistema deberá responder en menos de 2 segundos, sin tener en cuenta el tiempo que tarde la red en enviar la información de vuelta al usuario. \\ \hline
\centering {\bf Importancia} & Vital \\ \hline
\centering {\bf Urgencia} & Inmediatamente \\ \hline
\end{tabular}
\caption{Requisito no funcional 01. Tiempo de respuesta $<$ 2 segundos.}
\label{enlaceNFR1}
\end{center}
\end{table}

\begin{table}[H]
\begin{center}
\begin{tabular}{|p{3cm}|p{10cm}|} \hline
\centering {\bf NFR-02} & Los datos sensibles del usuario se almacenan encriptados  \\ \hline\hline
\centering {\bf Versión} & 1.0 (05/06/2022) \\ \hline
\centering {\bf Descripción} & El sistema deberá encriptar los datos sensibles del usuario como las contraseñas en la base de datos de lex.gal. Se utilizará un algoritmo SHA-256 para encriptar. \\ \hline
\centering {\bf Importancia} & Vital \\ \hline
\centering {\bf Urgencia} & Inmediatamente \\ \hline
\end{tabular}
\caption{Requisito no funcional 02. Los datos sensibles del usuario se almacenan encriptados.}
\label{enlaceNFR2}
\end{center}
\end{table}

\begin{table}[H]
\begin{center}
\begin{tabular}{|p{3cm}|p{10cm}|} \hline
\centering {\bf NFR-03} & No permitir que un rol realice operaciones de otro  \\ \hline\hline
\centering {\bf Versión} & 1.0 (05/06/2022) \\ \hline
\centering {\bf Descripción} & El sistema deberá restringir aquellas operaciones exclusivas del administrador para el resto de usuarios. \\ \hline
\centering {\bf Importancia} & Vital \\ \hline
\centering {\bf Urgencia} & Inmediatamente \\ \hline
\end{tabular}
\caption{Requisito no funcional 03. No permitir que un rol realice operaciones de otro.}
\label{enlaceNFR3}
\end{center}
\end{table}

\begin{table}[H]
\begin{center}
\begin{tabular}{|p{3cm}|p{10cm}|} \hline
\centering {\bf NFR-04} & Permanecer operativo durante el horario de trabajo \\ \hline\hline
\centering {\bf Versión} & 1.0 (05/06/2022) \\ \hline
\centering {\bf Descripción} & El sistema deberá permanecer operativo durante todo el horario de trabajo de los empleados de lex.gal.  \\ \hline
\centering {\bf Importancia} & Vital \\ \hline
\centering {\bf Urgencia} & Inmediatamente \\ \hline
\end{tabular}
\caption{Requisito no funcional 04. Permanecer operativo durante el horario de trabajo.}
\label{enlaceNFR4}
\end{center}
\end{table}

\begin{table}[H]
\begin{center}
\begin{tabular}{|p{3cm}|p{10cm}|} \hline
\centering {\bf NFR-05} & Tratamiento de documentos en lenguajes de etiquetas  \\ \hline\hline
\centering {\bf Versión} & 1.0 (05/06/2022) \\ \hline
\centering {\bf Descripción} & El sistema deberá restringir el tratamiento de documentos a documentos escritos en lenguajes de etiquetas, como HTML o XML. \\ \hline
\centering {\bf Importancia} & Vital \\ \hline
\centering {\bf Urgencia} & Inmediatamente \\ \hline
\end{tabular}
\caption{Requisito no funcional 05. Tratamiento de documentos en lenguajes de etiquetas.}
\label{enlaceNFR5}
\end{center}
\end{table}

\begin{table}[H]
\begin{center}
\begin{tabular}{|p{3cm}|p{10cm}|} \hline
\centering {\bf NFR-06} & Tiempo de aprendizaje por el usuario $<$ 1h  \\ \hline\hline
\centering {\bf Versión} & 1.0 (05/06/2022) \\ \hline
\centering {\bf Descripción} & El sistema deberá tener un proceso de aprendizaje rápido y sencillo para el usuario, que en todo caso deberá ser menor de 1 hora. \\ \hline
\centering {\bf Importancia} & Vital \\ \hline
\centering {\bf Urgencia} & Inmediatamente \\ \hline
\end{tabular}
\caption{Requisito no funcional 06. Tiempo de aprendizaje por el usuario $<$ 1h.}
\label{enlaceNFR6}
\end{center}
\end{table}

\begin{table}[H]
\begin{center}
\begin{tabular}{|p{3cm}|p{10cm}|} \hline
\centering {\bf NFR-07} & Proporcionar mensajes de apoyo al usuario  \\ \hline\hline
\centering {\bf Versión} & 1.0 (05/06/2022) \\ \hline
\centering {\bf Descripción} & El sistema deberá proporcionar mensajes de apoyo al usuario en todas las operaciones, con el objetivo de facilitar la navegación en la aplicación. \\ \hline
\centering {\bf Importancia} & Quedaría bien \\ \hline
\centering {\bf Urgencia} & Puede esperar \\ \hline
\end{tabular}
\caption{Requisito no funcional 07. Proporcionar mensajes de apoyo al usuario.}
\label{enlaceNFR7}
\end{center}
\end{table}

\begin{table}[H]
\begin{center}
\begin{tabular}{|p{3cm}|p{10cm}|} \hline
\centering {\bf NFR-08} & Uso de tecnologías web y servicios RESTful  \\ \hline\hline
\centering {\bf Versión} & 1.0 (05/06/2022) \\ \hline
\centering {\bf Descripción} & El sistema deberá emplear tecnologías web para realizar la aplicación, incluyendo HTML5, Tailwind como framework de CSS y React.js como framework de JavaScript. Además, todos los servicios serán RESTful, y serán implementados con Spring Boot. \\ \hline
\centering {\bf Importancia} & Vital \\ \hline
\centering {\bf Urgencia} & Inmediatamente \\ \hline
\end{tabular}
\caption{Requisito no funcional 08. Uso de tecnologías web y servicios RESTful.}
\label{enlaceNFR8}
\end{center}
\end{table}

\begin{table}[H]
\begin{center}
\begin{tabular}{|p{3cm}|p{10cm}|} \hline
\centering {\bf NFR-09} & Base de datos NoSQL  \\ \hline\hline
\centering {\bf Versión} & 1.0 (05/06/2022) \\ \hline
\centering {\bf Descripción} & El sistema deberá almacenar todos los datos en una base de datos NoSQL de MongoDB. \\ \hline
\centering {\bf Importancia} & Vital \\ \hline
\centering {\bf Urgencia} & Inmediatamente \\ \hline
\end{tabular}
\caption{Requisito no funcional 09. Base de datos NoSQL.}
\label{enlaceNFR9}
\end{center}
\end{table}

\begin{table}[H]
\begin{center}
\begin{tabular}{|p{3cm}|p{10cm}|} \hline
\centering {\bf NFR-10} & Generación de tokens JWT  \\ \hline\hline
\centering {\bf Versión} & 1.0 (05/06/2022) \\ \hline
\centering {\bf Descripción} & El sistema deberá generar tokens JWT cada vez que un usuario inicie sesión en la aplicación. De esta forma, los datos del usuario no quedan tan accesibles en la web, y será necesario utilizar un algoritmo  {\it hash} para desencriptarlos. En todo caso la contraseña no será almacenada en el token JWT. \\ \hline
\centering {\bf Importancia} & Vital \\ \hline
\centering {\bf Urgencia} & Inmediatamente \\ \hline
\end{tabular}
\caption{Requisito no funcional 10. Generación de tokens JWT.}
\label{enlaceNFR10}
\end{center}
\end{table}
\section{Requisitos de información}

Por último, los requisitos de información \cite{requisitos}, son aquellos requisitos que indican algún tipo de información, la cual es necesaria para el correcto funcionamiento del sistema.
\\

En el caso de esta aplicación, se han localizado dos: Usuario y Documento.
\\

Los requisitos de información pueden ser consultados en el \hyperref[APRequisitosInformacion]{Apéndice C.3. Requisitos de información}.
\section{Casos de uso}

Un caso de uso \cite{casodeuso} es cualquier acción o actividad que se realiza en el sistema. Representan diferentes operaciones que se realizan en la aplicación.
\\

Los casos de uso pueden ser consultados en el \hyperref[APCasosUso]{Apéndice C.4. Casos de uso}.
\cleardoublepage
\chapter{Especificación de requisitos}
\label{enlaceespecificacion}

Para proceder a especificar los requisitos se ha realizado un profundo análisis de los objetivos de la aplicación. Estos requisitos se han buscado teniendo en cuenta aspectos funcionales, de rendimiento, de seguridad, de información, etc.
\\

Las plantillas empleadas para las tablas de los requisitos y casos de uso se encuentran en los apéndices: la \hyperref[enlaceFRQX]{E.1. Plantilla de requisitos funcionales}, la \hyperref[enlaceNFRX]{E.2. Plantilla de requisitos no funcionales}, la \hyperref[enlaceIRQX]{E.3. Plantilla de requisitos de información} y la \hyperref[enlaceUCX]{E.4. Plantilla de casos de uso}.
Estas plantillas están inspiradas en las tablas empleadas en el REM \cite{rem} para la gestión de requisitos.
\\

La importancia de los requisitos será medida según el grado de la necesidad del requisito en la aplicación, siendo en el caso de este proyecto ``Vital" para los requisitos de mayor importancia, y, ``Quedaría bien" para los menos importantes. Por su parte, la urgencia, indicará aquellas partes que se han de implementar con mayor rapidez, siendo ``Inmediatamente" las que tienen una mayor prioridad, y, ``Puede esperar", para los requisitos que no son estrictamente necesarios.

\section{Requisitos funcionales}

Los requisitos funcionales \cite{requisitos} son aquellos que definen una funcionalidad del sistema de software o sus componentes, tanto para las entradas de datos como para el comportamiento en situaciones específicas.
\\

En el caso de esta aplicación se han localizado dos grupos de requisitos funcionales: los asociados con operaciones de usuarios, y los que se relacionan con operaciones asociadas a leyes.
\\

Los requisitos funcionales pueden ser consultados en el \hyperref[APRequisitosFuncionales]{Apéndice C.1. Requisitos funcionales}.
\section{Requisitos no funcionales}
\label{APRequisitosNoFuncionales}

\begin{table}[H]
\begin{center}
\begin{tabular}{|p{3cm}|p{10cm}|} \hline
\centering {\bf NFR-01} & Tiempo de respuesta $<$ 2 segundos  \\ \hline\hline
\centering {\bf Versión} & 1.0 (05/06/2022) \\ \hline
\centering {\bf Descripción} & El sistema deberá responder en menos de 2 segundos, sin tener en cuenta el tiempo que tarde la red en enviar la información de vuelta al usuario. \\ \hline
\centering {\bf Importancia} & Vital \\ \hline
\centering {\bf Urgencia} & Inmediatamente \\ \hline
\end{tabular}
\caption{Requisito no funcional 01. Tiempo de respuesta $<$ 2 segundos.}
\label{enlaceNFR1}
\end{center}
\end{table}

\begin{table}[H]
\begin{center}
\begin{tabular}{|p{3cm}|p{10cm}|} \hline
\centering {\bf NFR-02} & Los datos sensibles del usuario se almacenan encriptados  \\ \hline\hline
\centering {\bf Versión} & 1.0 (05/06/2022) \\ \hline
\centering {\bf Descripción} & El sistema deberá encriptar los datos sensibles del usuario como las contraseñas en la base de datos de lex.gal. Se utilizará un algoritmo SHA-256 para encriptar. \\ \hline
\centering {\bf Importancia} & Vital \\ \hline
\centering {\bf Urgencia} & Inmediatamente \\ \hline
\end{tabular}
\caption{Requisito no funcional 02. Los datos sensibles del usuario se almacenan encriptados.}
\label{enlaceNFR2}
\end{center}
\end{table}

\begin{table}[H]
\begin{center}
\begin{tabular}{|p{3cm}|p{10cm}|} \hline
\centering {\bf NFR-03} & No permitir que un rol realice operaciones de otro  \\ \hline\hline
\centering {\bf Versión} & 1.0 (05/06/2022) \\ \hline
\centering {\bf Descripción} & El sistema deberá restringir aquellas operaciones exclusivas del administrador para el resto de usuarios. \\ \hline
\centering {\bf Importancia} & Vital \\ \hline
\centering {\bf Urgencia} & Inmediatamente \\ \hline
\end{tabular}
\caption{Requisito no funcional 03. No permitir que un rol realice operaciones de otro.}
\label{enlaceNFR3}
\end{center}
\end{table}

\begin{table}[H]
\begin{center}
\begin{tabular}{|p{3cm}|p{10cm}|} \hline
\centering {\bf NFR-04} & Permanecer operativo durante el horario de trabajo \\ \hline\hline
\centering {\bf Versión} & 1.0 (05/06/2022) \\ \hline
\centering {\bf Descripción} & El sistema deberá permanecer operativo durante todo el horario de trabajo de los empleados de lex.gal.  \\ \hline
\centering {\bf Importancia} & Vital \\ \hline
\centering {\bf Urgencia} & Inmediatamente \\ \hline
\end{tabular}
\caption{Requisito no funcional 04. Permanecer operativo durante el horario de trabajo.}
\label{enlaceNFR4}
\end{center}
\end{table}

\begin{table}[H]
\begin{center}
\begin{tabular}{|p{3cm}|p{10cm}|} \hline
\centering {\bf NFR-05} & Tratamiento de documentos en lenguajes de etiquetas  \\ \hline\hline
\centering {\bf Versión} & 1.0 (05/06/2022) \\ \hline
\centering {\bf Descripción} & El sistema deberá restringir el tratamiento de documentos a documentos escritos en lenguajes de etiquetas, como HTML o XML. \\ \hline
\centering {\bf Importancia} & Vital \\ \hline
\centering {\bf Urgencia} & Inmediatamente \\ \hline
\end{tabular}
\caption{Requisito no funcional 05. Tratamiento de documentos en lenguajes de etiquetas.}
\label{enlaceNFR5}
\end{center}
\end{table}

\begin{table}[H]
\begin{center}
\begin{tabular}{|p{3cm}|p{10cm}|} \hline
\centering {\bf NFR-06} & Tiempo de aprendizaje por el usuario $<$ 1h  \\ \hline\hline
\centering {\bf Versión} & 1.0 (05/06/2022) \\ \hline
\centering {\bf Descripción} & El sistema deberá tener un proceso de aprendizaje rápido y sencillo para el usuario, que en todo caso deberá ser menor de 1 hora. \\ \hline
\centering {\bf Importancia} & Vital \\ \hline
\centering {\bf Urgencia} & Inmediatamente \\ \hline
\end{tabular}
\caption{Requisito no funcional 06. Tiempo de aprendizaje por el usuario $<$ 1h.}
\label{enlaceNFR6}
\end{center}
\end{table}

\begin{table}[H]
\begin{center}
\begin{tabular}{|p{3cm}|p{10cm}|} \hline
\centering {\bf NFR-07} & Proporcionar mensajes de apoyo al usuario  \\ \hline\hline
\centering {\bf Versión} & 1.0 (05/06/2022) \\ \hline
\centering {\bf Descripción} & El sistema deberá proporcionar mensajes de apoyo al usuario en todas las operaciones, con el objetivo de facilitar la navegación en la aplicación. \\ \hline
\centering {\bf Importancia} & Quedaría bien \\ \hline
\centering {\bf Urgencia} & Puede esperar \\ \hline
\end{tabular}
\caption{Requisito no funcional 07. Proporcionar mensajes de apoyo al usuario.}
\label{enlaceNFR7}
\end{center}
\end{table}

\begin{table}[H]
\begin{center}
\begin{tabular}{|p{3cm}|p{10cm}|} \hline
\centering {\bf NFR-08} & Uso de tecnologías web y servicios RESTful  \\ \hline\hline
\centering {\bf Versión} & 1.0 (05/06/2022) \\ \hline
\centering {\bf Descripción} & El sistema deberá emplear tecnologías web para realizar la aplicación, incluyendo HTML5, Tailwind como framework de CSS y React.js como framework de JavaScript. Además, todos los servicios serán RESTful, y serán implementados con Spring Boot. \\ \hline
\centering {\bf Importancia} & Vital \\ \hline
\centering {\bf Urgencia} & Inmediatamente \\ \hline
\end{tabular}
\caption{Requisito no funcional 08. Uso de tecnologías web y servicios RESTful.}
\label{enlaceNFR8}
\end{center}
\end{table}

\begin{table}[H]
\begin{center}
\begin{tabular}{|p{3cm}|p{10cm}|} \hline
\centering {\bf NFR-09} & Base de datos NoSQL  \\ \hline\hline
\centering {\bf Versión} & 1.0 (05/06/2022) \\ \hline
\centering {\bf Descripción} & El sistema deberá almacenar todos los datos en una base de datos NoSQL de MongoDB. \\ \hline
\centering {\bf Importancia} & Vital \\ \hline
\centering {\bf Urgencia} & Inmediatamente \\ \hline
\end{tabular}
\caption{Requisito no funcional 09. Base de datos NoSQL.}
\label{enlaceNFR9}
\end{center}
\end{table}

\begin{table}[H]
\begin{center}
\begin{tabular}{|p{3cm}|p{10cm}|} \hline
\centering {\bf NFR-10} & Generación de tokens JWT  \\ \hline\hline
\centering {\bf Versión} & 1.0 (05/06/2022) \\ \hline
\centering {\bf Descripción} & El sistema deberá generar tokens JWT cada vez que un usuario inicie sesión en la aplicación. De esta forma, los datos del usuario no quedan tan accesibles en la web, y será necesario utilizar un algoritmo  {\it hash} para desencriptarlos. En todo caso la contraseña no será almacenada en el token JWT. \\ \hline
\centering {\bf Importancia} & Vital \\ \hline
\centering {\bf Urgencia} & Inmediatamente \\ \hline
\end{tabular}
\caption{Requisito no funcional 10. Generación de tokens JWT.}
\label{enlaceNFR10}
\end{center}
\end{table}
\section{Requisitos de información}

Por último, los requisitos de información \cite{requisitos}, son aquellos requisitos que indican algún tipo de información, la cual es necesaria para el correcto funcionamiento del sistema.
\\

En el caso de esta aplicación, se han localizado dos: Usuario y Documento.
\\

Los requisitos de información pueden ser consultados en el \hyperref[APRequisitosInformacion]{Apéndice C.3. Requisitos de información}.
\section{Casos de uso}

Un caso de uso \cite{casodeuso} es cualquier acción o actividad que se realiza en el sistema. Representan diferentes operaciones que se realizan en la aplicación.
\\

Los casos de uso pueden ser consultados en el \hyperref[APCasosUso]{Apéndice C.4. Casos de uso}.
\cleardoublepage
\chapter{Especificación de requisitos}
\label{enlaceespecificacion}

Para proceder a especificar los requisitos se ha realizado un profundo análisis de los objetivos de la aplicación. Estos requisitos se han buscado teniendo en cuenta aspectos funcionales, de rendimiento, de seguridad, de información, etc.
\\

Las plantillas empleadas para las tablas de los requisitos y casos de uso se encuentran en los apéndices: la \hyperref[enlaceFRQX]{E.1. Plantilla de requisitos funcionales}, la \hyperref[enlaceNFRX]{E.2. Plantilla de requisitos no funcionales}, la \hyperref[enlaceIRQX]{E.3. Plantilla de requisitos de información} y la \hyperref[enlaceUCX]{E.4. Plantilla de casos de uso}.
Estas plantillas están inspiradas en las tablas empleadas en el REM \cite{rem} para la gestión de requisitos.
\\

La importancia de los requisitos será medida según el grado de la necesidad del requisito en la aplicación, siendo en el caso de este proyecto ``Vital" para los requisitos de mayor importancia, y, ``Quedaría bien" para los menos importantes. Por su parte, la urgencia, indicará aquellas partes que se han de implementar con mayor rapidez, siendo ``Inmediatamente" las que tienen una mayor prioridad, y, ``Puede esperar", para los requisitos que no son estrictamente necesarios.

\section{Requisitos funcionales}

Los requisitos funcionales \cite{requisitos} son aquellos que definen una funcionalidad del sistema de software o sus componentes, tanto para las entradas de datos como para el comportamiento en situaciones específicas.
\\

En el caso de esta aplicación se han localizado dos grupos de requisitos funcionales: los asociados con operaciones de usuarios, y los que se relacionan con operaciones asociadas a leyes.
\\

Los requisitos funcionales pueden ser consultados en el \hyperref[APRequisitosFuncionales]{Apéndice C.1. Requisitos funcionales}.
\section{Requisitos no funcionales}
\label{APRequisitosNoFuncionales}

\begin{table}[H]
\begin{center}
\begin{tabular}{|p{3cm}|p{10cm}|} \hline
\centering {\bf NFR-01} & Tiempo de respuesta $<$ 2 segundos  \\ \hline\hline
\centering {\bf Versión} & 1.0 (05/06/2022) \\ \hline
\centering {\bf Descripción} & El sistema deberá responder en menos de 2 segundos, sin tener en cuenta el tiempo que tarde la red en enviar la información de vuelta al usuario. \\ \hline
\centering {\bf Importancia} & Vital \\ \hline
\centering {\bf Urgencia} & Inmediatamente \\ \hline
\end{tabular}
\caption{Requisito no funcional 01. Tiempo de respuesta $<$ 2 segundos.}
\label{enlaceNFR1}
\end{center}
\end{table}

\begin{table}[H]
\begin{center}
\begin{tabular}{|p{3cm}|p{10cm}|} \hline
\centering {\bf NFR-02} & Los datos sensibles del usuario se almacenan encriptados  \\ \hline\hline
\centering {\bf Versión} & 1.0 (05/06/2022) \\ \hline
\centering {\bf Descripción} & El sistema deberá encriptar los datos sensibles del usuario como las contraseñas en la base de datos de lex.gal. Se utilizará un algoritmo SHA-256 para encriptar. \\ \hline
\centering {\bf Importancia} & Vital \\ \hline
\centering {\bf Urgencia} & Inmediatamente \\ \hline
\end{tabular}
\caption{Requisito no funcional 02. Los datos sensibles del usuario se almacenan encriptados.}
\label{enlaceNFR2}
\end{center}
\end{table}

\begin{table}[H]
\begin{center}
\begin{tabular}{|p{3cm}|p{10cm}|} \hline
\centering {\bf NFR-03} & No permitir que un rol realice operaciones de otro  \\ \hline\hline
\centering {\bf Versión} & 1.0 (05/06/2022) \\ \hline
\centering {\bf Descripción} & El sistema deberá restringir aquellas operaciones exclusivas del administrador para el resto de usuarios. \\ \hline
\centering {\bf Importancia} & Vital \\ \hline
\centering {\bf Urgencia} & Inmediatamente \\ \hline
\end{tabular}
\caption{Requisito no funcional 03. No permitir que un rol realice operaciones de otro.}
\label{enlaceNFR3}
\end{center}
\end{table}

\begin{table}[H]
\begin{center}
\begin{tabular}{|p{3cm}|p{10cm}|} \hline
\centering {\bf NFR-04} & Permanecer operativo durante el horario de trabajo \\ \hline\hline
\centering {\bf Versión} & 1.0 (05/06/2022) \\ \hline
\centering {\bf Descripción} & El sistema deberá permanecer operativo durante todo el horario de trabajo de los empleados de lex.gal.  \\ \hline
\centering {\bf Importancia} & Vital \\ \hline
\centering {\bf Urgencia} & Inmediatamente \\ \hline
\end{tabular}
\caption{Requisito no funcional 04. Permanecer operativo durante el horario de trabajo.}
\label{enlaceNFR4}
\end{center}
\end{table}

\begin{table}[H]
\begin{center}
\begin{tabular}{|p{3cm}|p{10cm}|} \hline
\centering {\bf NFR-05} & Tratamiento de documentos en lenguajes de etiquetas  \\ \hline\hline
\centering {\bf Versión} & 1.0 (05/06/2022) \\ \hline
\centering {\bf Descripción} & El sistema deberá restringir el tratamiento de documentos a documentos escritos en lenguajes de etiquetas, como HTML o XML. \\ \hline
\centering {\bf Importancia} & Vital \\ \hline
\centering {\bf Urgencia} & Inmediatamente \\ \hline
\end{tabular}
\caption{Requisito no funcional 05. Tratamiento de documentos en lenguajes de etiquetas.}
\label{enlaceNFR5}
\end{center}
\end{table}

\begin{table}[H]
\begin{center}
\begin{tabular}{|p{3cm}|p{10cm}|} \hline
\centering {\bf NFR-06} & Tiempo de aprendizaje por el usuario $<$ 1h  \\ \hline\hline
\centering {\bf Versión} & 1.0 (05/06/2022) \\ \hline
\centering {\bf Descripción} & El sistema deberá tener un proceso de aprendizaje rápido y sencillo para el usuario, que en todo caso deberá ser menor de 1 hora. \\ \hline
\centering {\bf Importancia} & Vital \\ \hline
\centering {\bf Urgencia} & Inmediatamente \\ \hline
\end{tabular}
\caption{Requisito no funcional 06. Tiempo de aprendizaje por el usuario $<$ 1h.}
\label{enlaceNFR6}
\end{center}
\end{table}

\begin{table}[H]
\begin{center}
\begin{tabular}{|p{3cm}|p{10cm}|} \hline
\centering {\bf NFR-07} & Proporcionar mensajes de apoyo al usuario  \\ \hline\hline
\centering {\bf Versión} & 1.0 (05/06/2022) \\ \hline
\centering {\bf Descripción} & El sistema deberá proporcionar mensajes de apoyo al usuario en todas las operaciones, con el objetivo de facilitar la navegación en la aplicación. \\ \hline
\centering {\bf Importancia} & Quedaría bien \\ \hline
\centering {\bf Urgencia} & Puede esperar \\ \hline
\end{tabular}
\caption{Requisito no funcional 07. Proporcionar mensajes de apoyo al usuario.}
\label{enlaceNFR7}
\end{center}
\end{table}

\begin{table}[H]
\begin{center}
\begin{tabular}{|p{3cm}|p{10cm}|} \hline
\centering {\bf NFR-08} & Uso de tecnologías web y servicios RESTful  \\ \hline\hline
\centering {\bf Versión} & 1.0 (05/06/2022) \\ \hline
\centering {\bf Descripción} & El sistema deberá emplear tecnologías web para realizar la aplicación, incluyendo HTML5, Tailwind como framework de CSS y React.js como framework de JavaScript. Además, todos los servicios serán RESTful, y serán implementados con Spring Boot. \\ \hline
\centering {\bf Importancia} & Vital \\ \hline
\centering {\bf Urgencia} & Inmediatamente \\ \hline
\end{tabular}
\caption{Requisito no funcional 08. Uso de tecnologías web y servicios RESTful.}
\label{enlaceNFR8}
\end{center}
\end{table}

\begin{table}[H]
\begin{center}
\begin{tabular}{|p{3cm}|p{10cm}|} \hline
\centering {\bf NFR-09} & Base de datos NoSQL  \\ \hline\hline
\centering {\bf Versión} & 1.0 (05/06/2022) \\ \hline
\centering {\bf Descripción} & El sistema deberá almacenar todos los datos en una base de datos NoSQL de MongoDB. \\ \hline
\centering {\bf Importancia} & Vital \\ \hline
\centering {\bf Urgencia} & Inmediatamente \\ \hline
\end{tabular}
\caption{Requisito no funcional 09. Base de datos NoSQL.}
\label{enlaceNFR9}
\end{center}
\end{table}

\begin{table}[H]
\begin{center}
\begin{tabular}{|p{3cm}|p{10cm}|} \hline
\centering {\bf NFR-10} & Generación de tokens JWT  \\ \hline\hline
\centering {\bf Versión} & 1.0 (05/06/2022) \\ \hline
\centering {\bf Descripción} & El sistema deberá generar tokens JWT cada vez que un usuario inicie sesión en la aplicación. De esta forma, los datos del usuario no quedan tan accesibles en la web, y será necesario utilizar un algoritmo  {\it hash} para desencriptarlos. En todo caso la contraseña no será almacenada en el token JWT. \\ \hline
\centering {\bf Importancia} & Vital \\ \hline
\centering {\bf Urgencia} & Inmediatamente \\ \hline
\end{tabular}
\caption{Requisito no funcional 10. Generación de tokens JWT.}
\label{enlaceNFR10}
\end{center}
\end{table}
\section{Requisitos de información}

Por último, los requisitos de información \cite{requisitos}, son aquellos requisitos que indican algún tipo de información, la cual es necesaria para el correcto funcionamiento del sistema.
\\

En el caso de esta aplicación, se han localizado dos: Usuario y Documento.
\\

Los requisitos de información pueden ser consultados en el \hyperref[APRequisitosInformacion]{Apéndice C.3. Requisitos de información}.
\section{Casos de uso}

Un caso de uso \cite{casodeuso} es cualquier acción o actividad que se realiza en el sistema. Representan diferentes operaciones que se realizan en la aplicación.
\\

Los casos de uso pueden ser consultados en el \hyperref[APCasosUso]{Apéndice C.4. Casos de uso}.
\cleardoublepage
\chapter{Especificación de requisitos}
\label{enlaceespecificacion}

Para proceder a especificar los requisitos se ha realizado un profundo análisis de los objetivos de la aplicación. Estos requisitos se han buscado teniendo en cuenta aspectos funcionales, de rendimiento, de seguridad, de información, etc.
\\

Las plantillas empleadas para las tablas de los requisitos y casos de uso se encuentran en los apéndices: la \hyperref[enlaceFRQX]{E.1. Plantilla de requisitos funcionales}, la \hyperref[enlaceNFRX]{E.2. Plantilla de requisitos no funcionales}, la \hyperref[enlaceIRQX]{E.3. Plantilla de requisitos de información} y la \hyperref[enlaceUCX]{E.4. Plantilla de casos de uso}.
Estas plantillas están inspiradas en las tablas empleadas en el REM \cite{rem} para la gestión de requisitos.
\\

La importancia de los requisitos será medida según el grado de la necesidad del requisito en la aplicación, siendo en el caso de este proyecto ``Vital" para los requisitos de mayor importancia, y, ``Quedaría bien" para los menos importantes. Por su parte, la urgencia, indicará aquellas partes que se han de implementar con mayor rapidez, siendo ``Inmediatamente" las que tienen una mayor prioridad, y, ``Puede esperar", para los requisitos que no son estrictamente necesarios.

\section{Requisitos funcionales}

Los requisitos funcionales \cite{requisitos} son aquellos que definen una funcionalidad del sistema de software o sus componentes, tanto para las entradas de datos como para el comportamiento en situaciones específicas.
\\

En el caso de esta aplicación se han localizado dos grupos de requisitos funcionales: los asociados con operaciones de usuarios, y los que se relacionan con operaciones asociadas a leyes.
\\

Los requisitos funcionales pueden ser consultados en el \hyperref[APRequisitosFuncionales]{Apéndice C.1. Requisitos funcionales}.
\section{Requisitos no funcionales}
\label{APRequisitosNoFuncionales}

\begin{table}[H]
\begin{center}
\begin{tabular}{|p{3cm}|p{10cm}|} \hline
\centering {\bf NFR-01} & Tiempo de respuesta $<$ 2 segundos  \\ \hline\hline
\centering {\bf Versión} & 1.0 (05/06/2022) \\ \hline
\centering {\bf Descripción} & El sistema deberá responder en menos de 2 segundos, sin tener en cuenta el tiempo que tarde la red en enviar la información de vuelta al usuario. \\ \hline
\centering {\bf Importancia} & Vital \\ \hline
\centering {\bf Urgencia} & Inmediatamente \\ \hline
\end{tabular}
\caption{Requisito no funcional 01. Tiempo de respuesta $<$ 2 segundos.}
\label{enlaceNFR1}
\end{center}
\end{table}

\begin{table}[H]
\begin{center}
\begin{tabular}{|p{3cm}|p{10cm}|} \hline
\centering {\bf NFR-02} & Los datos sensibles del usuario se almacenan encriptados  \\ \hline\hline
\centering {\bf Versión} & 1.0 (05/06/2022) \\ \hline
\centering {\bf Descripción} & El sistema deberá encriptar los datos sensibles del usuario como las contraseñas en la base de datos de lex.gal. Se utilizará un algoritmo SHA-256 para encriptar. \\ \hline
\centering {\bf Importancia} & Vital \\ \hline
\centering {\bf Urgencia} & Inmediatamente \\ \hline
\end{tabular}
\caption{Requisito no funcional 02. Los datos sensibles del usuario se almacenan encriptados.}
\label{enlaceNFR2}
\end{center}
\end{table}

\begin{table}[H]
\begin{center}
\begin{tabular}{|p{3cm}|p{10cm}|} \hline
\centering {\bf NFR-03} & No permitir que un rol realice operaciones de otro  \\ \hline\hline
\centering {\bf Versión} & 1.0 (05/06/2022) \\ \hline
\centering {\bf Descripción} & El sistema deberá restringir aquellas operaciones exclusivas del administrador para el resto de usuarios. \\ \hline
\centering {\bf Importancia} & Vital \\ \hline
\centering {\bf Urgencia} & Inmediatamente \\ \hline
\end{tabular}
\caption{Requisito no funcional 03. No permitir que un rol realice operaciones de otro.}
\label{enlaceNFR3}
\end{center}
\end{table}

\begin{table}[H]
\begin{center}
\begin{tabular}{|p{3cm}|p{10cm}|} \hline
\centering {\bf NFR-04} & Permanecer operativo durante el horario de trabajo \\ \hline\hline
\centering {\bf Versión} & 1.0 (05/06/2022) \\ \hline
\centering {\bf Descripción} & El sistema deberá permanecer operativo durante todo el horario de trabajo de los empleados de lex.gal.  \\ \hline
\centering {\bf Importancia} & Vital \\ \hline
\centering {\bf Urgencia} & Inmediatamente \\ \hline
\end{tabular}
\caption{Requisito no funcional 04. Permanecer operativo durante el horario de trabajo.}
\label{enlaceNFR4}
\end{center}
\end{table}

\begin{table}[H]
\begin{center}
\begin{tabular}{|p{3cm}|p{10cm}|} \hline
\centering {\bf NFR-05} & Tratamiento de documentos en lenguajes de etiquetas  \\ \hline\hline
\centering {\bf Versión} & 1.0 (05/06/2022) \\ \hline
\centering {\bf Descripción} & El sistema deberá restringir el tratamiento de documentos a documentos escritos en lenguajes de etiquetas, como HTML o XML. \\ \hline
\centering {\bf Importancia} & Vital \\ \hline
\centering {\bf Urgencia} & Inmediatamente \\ \hline
\end{tabular}
\caption{Requisito no funcional 05. Tratamiento de documentos en lenguajes de etiquetas.}
\label{enlaceNFR5}
\end{center}
\end{table}

\begin{table}[H]
\begin{center}
\begin{tabular}{|p{3cm}|p{10cm}|} \hline
\centering {\bf NFR-06} & Tiempo de aprendizaje por el usuario $<$ 1h  \\ \hline\hline
\centering {\bf Versión} & 1.0 (05/06/2022) \\ \hline
\centering {\bf Descripción} & El sistema deberá tener un proceso de aprendizaje rápido y sencillo para el usuario, que en todo caso deberá ser menor de 1 hora. \\ \hline
\centering {\bf Importancia} & Vital \\ \hline
\centering {\bf Urgencia} & Inmediatamente \\ \hline
\end{tabular}
\caption{Requisito no funcional 06. Tiempo de aprendizaje por el usuario $<$ 1h.}
\label{enlaceNFR6}
\end{center}
\end{table}

\begin{table}[H]
\begin{center}
\begin{tabular}{|p{3cm}|p{10cm}|} \hline
\centering {\bf NFR-07} & Proporcionar mensajes de apoyo al usuario  \\ \hline\hline
\centering {\bf Versión} & 1.0 (05/06/2022) \\ \hline
\centering {\bf Descripción} & El sistema deberá proporcionar mensajes de apoyo al usuario en todas las operaciones, con el objetivo de facilitar la navegación en la aplicación. \\ \hline
\centering {\bf Importancia} & Quedaría bien \\ \hline
\centering {\bf Urgencia} & Puede esperar \\ \hline
\end{tabular}
\caption{Requisito no funcional 07. Proporcionar mensajes de apoyo al usuario.}
\label{enlaceNFR7}
\end{center}
\end{table}

\begin{table}[H]
\begin{center}
\begin{tabular}{|p{3cm}|p{10cm}|} \hline
\centering {\bf NFR-08} & Uso de tecnologías web y servicios RESTful  \\ \hline\hline
\centering {\bf Versión} & 1.0 (05/06/2022) \\ \hline
\centering {\bf Descripción} & El sistema deberá emplear tecnologías web para realizar la aplicación, incluyendo HTML5, Tailwind como framework de CSS y React.js como framework de JavaScript. Además, todos los servicios serán RESTful, y serán implementados con Spring Boot. \\ \hline
\centering {\bf Importancia} & Vital \\ \hline
\centering {\bf Urgencia} & Inmediatamente \\ \hline
\end{tabular}
\caption{Requisito no funcional 08. Uso de tecnologías web y servicios RESTful.}
\label{enlaceNFR8}
\end{center}
\end{table}

\begin{table}[H]
\begin{center}
\begin{tabular}{|p{3cm}|p{10cm}|} \hline
\centering {\bf NFR-09} & Base de datos NoSQL  \\ \hline\hline
\centering {\bf Versión} & 1.0 (05/06/2022) \\ \hline
\centering {\bf Descripción} & El sistema deberá almacenar todos los datos en una base de datos NoSQL de MongoDB. \\ \hline
\centering {\bf Importancia} & Vital \\ \hline
\centering {\bf Urgencia} & Inmediatamente \\ \hline
\end{tabular}
\caption{Requisito no funcional 09. Base de datos NoSQL.}
\label{enlaceNFR9}
\end{center}
\end{table}

\begin{table}[H]
\begin{center}
\begin{tabular}{|p{3cm}|p{10cm}|} \hline
\centering {\bf NFR-10} & Generación de tokens JWT  \\ \hline\hline
\centering {\bf Versión} & 1.0 (05/06/2022) \\ \hline
\centering {\bf Descripción} & El sistema deberá generar tokens JWT cada vez que un usuario inicie sesión en la aplicación. De esta forma, los datos del usuario no quedan tan accesibles en la web, y será necesario utilizar un algoritmo  {\it hash} para desencriptarlos. En todo caso la contraseña no será almacenada en el token JWT. \\ \hline
\centering {\bf Importancia} & Vital \\ \hline
\centering {\bf Urgencia} & Inmediatamente \\ \hline
\end{tabular}
\caption{Requisito no funcional 10. Generación de tokens JWT.}
\label{enlaceNFR10}
\end{center}
\end{table}
\section{Requisitos de información}

Por último, los requisitos de información \cite{requisitos}, son aquellos requisitos que indican algún tipo de información, la cual es necesaria para el correcto funcionamiento del sistema.
\\

En el caso de esta aplicación, se han localizado dos: Usuario y Documento.
\\

Los requisitos de información pueden ser consultados en el \hyperref[APRequisitosInformacion]{Apéndice C.3. Requisitos de información}.
\section{Casos de uso}

Un caso de uso \cite{casodeuso} es cualquier acción o actividad que se realiza en el sistema. Representan diferentes operaciones que se realizan en la aplicación.
\\

Los casos de uso pueden ser consultados en el \hyperref[APCasosUso]{Apéndice C.4. Casos de uso}.
\cleardoublepage
\chapter{Especificación de requisitos}
\label{enlaceespecificacion}

Para proceder a especificar los requisitos se ha realizado un profundo análisis de los objetivos de la aplicación. Estos requisitos se han buscado teniendo en cuenta aspectos funcionales, de rendimiento, de seguridad, de información, etc.
\\

Las plantillas empleadas para las tablas de los requisitos y casos de uso se encuentran en los apéndices: la \hyperref[enlaceFRQX]{E.1. Plantilla de requisitos funcionales}, la \hyperref[enlaceNFRX]{E.2. Plantilla de requisitos no funcionales}, la \hyperref[enlaceIRQX]{E.3. Plantilla de requisitos de información} y la \hyperref[enlaceUCX]{E.4. Plantilla de casos de uso}.
Estas plantillas están inspiradas en las tablas empleadas en el REM \cite{rem} para la gestión de requisitos.
\\

La importancia de los requisitos será medida según el grado de la necesidad del requisito en la aplicación, siendo en el caso de este proyecto ``Vital" para los requisitos de mayor importancia, y, ``Quedaría bien" para los menos importantes. Por su parte, la urgencia, indicará aquellas partes que se han de implementar con mayor rapidez, siendo ``Inmediatamente" las que tienen una mayor prioridad, y, ``Puede esperar", para los requisitos que no son estrictamente necesarios.

\section{Requisitos funcionales}

Los requisitos funcionales \cite{requisitos} son aquellos que definen una funcionalidad del sistema de software o sus componentes, tanto para las entradas de datos como para el comportamiento en situaciones específicas.
\\

En el caso de esta aplicación se han localizado dos grupos de requisitos funcionales: los asociados con operaciones de usuarios, y los que se relacionan con operaciones asociadas a leyes.
\\

Los requisitos funcionales pueden ser consultados en el \hyperref[APRequisitosFuncionales]{Apéndice C.1. Requisitos funcionales}.
\section{Requisitos no funcionales}
\label{APRequisitosNoFuncionales}

\begin{table}[H]
\begin{center}
\begin{tabular}{|p{3cm}|p{10cm}|} \hline
\centering {\bf NFR-01} & Tiempo de respuesta $<$ 2 segundos  \\ \hline\hline
\centering {\bf Versión} & 1.0 (05/06/2022) \\ \hline
\centering {\bf Descripción} & El sistema deberá responder en menos de 2 segundos, sin tener en cuenta el tiempo que tarde la red en enviar la información de vuelta al usuario. \\ \hline
\centering {\bf Importancia} & Vital \\ \hline
\centering {\bf Urgencia} & Inmediatamente \\ \hline
\end{tabular}
\caption{Requisito no funcional 01. Tiempo de respuesta $<$ 2 segundos.}
\label{enlaceNFR1}
\end{center}
\end{table}

\begin{table}[H]
\begin{center}
\begin{tabular}{|p{3cm}|p{10cm}|} \hline
\centering {\bf NFR-02} & Los datos sensibles del usuario se almacenan encriptados  \\ \hline\hline
\centering {\bf Versión} & 1.0 (05/06/2022) \\ \hline
\centering {\bf Descripción} & El sistema deberá encriptar los datos sensibles del usuario como las contraseñas en la base de datos de lex.gal. Se utilizará un algoritmo SHA-256 para encriptar. \\ \hline
\centering {\bf Importancia} & Vital \\ \hline
\centering {\bf Urgencia} & Inmediatamente \\ \hline
\end{tabular}
\caption{Requisito no funcional 02. Los datos sensibles del usuario se almacenan encriptados.}
\label{enlaceNFR2}
\end{center}
\end{table}

\begin{table}[H]
\begin{center}
\begin{tabular}{|p{3cm}|p{10cm}|} \hline
\centering {\bf NFR-03} & No permitir que un rol realice operaciones de otro  \\ \hline\hline
\centering {\bf Versión} & 1.0 (05/06/2022) \\ \hline
\centering {\bf Descripción} & El sistema deberá restringir aquellas operaciones exclusivas del administrador para el resto de usuarios. \\ \hline
\centering {\bf Importancia} & Vital \\ \hline
\centering {\bf Urgencia} & Inmediatamente \\ \hline
\end{tabular}
\caption{Requisito no funcional 03. No permitir que un rol realice operaciones de otro.}
\label{enlaceNFR3}
\end{center}
\end{table}

\begin{table}[H]
\begin{center}
\begin{tabular}{|p{3cm}|p{10cm}|} \hline
\centering {\bf NFR-04} & Permanecer operativo durante el horario de trabajo \\ \hline\hline
\centering {\bf Versión} & 1.0 (05/06/2022) \\ \hline
\centering {\bf Descripción} & El sistema deberá permanecer operativo durante todo el horario de trabajo de los empleados de lex.gal.  \\ \hline
\centering {\bf Importancia} & Vital \\ \hline
\centering {\bf Urgencia} & Inmediatamente \\ \hline
\end{tabular}
\caption{Requisito no funcional 04. Permanecer operativo durante el horario de trabajo.}
\label{enlaceNFR4}
\end{center}
\end{table}

\begin{table}[H]
\begin{center}
\begin{tabular}{|p{3cm}|p{10cm}|} \hline
\centering {\bf NFR-05} & Tratamiento de documentos en lenguajes de etiquetas  \\ \hline\hline
\centering {\bf Versión} & 1.0 (05/06/2022) \\ \hline
\centering {\bf Descripción} & El sistema deberá restringir el tratamiento de documentos a documentos escritos en lenguajes de etiquetas, como HTML o XML. \\ \hline
\centering {\bf Importancia} & Vital \\ \hline
\centering {\bf Urgencia} & Inmediatamente \\ \hline
\end{tabular}
\caption{Requisito no funcional 05. Tratamiento de documentos en lenguajes de etiquetas.}
\label{enlaceNFR5}
\end{center}
\end{table}

\begin{table}[H]
\begin{center}
\begin{tabular}{|p{3cm}|p{10cm}|} \hline
\centering {\bf NFR-06} & Tiempo de aprendizaje por el usuario $<$ 1h  \\ \hline\hline
\centering {\bf Versión} & 1.0 (05/06/2022) \\ \hline
\centering {\bf Descripción} & El sistema deberá tener un proceso de aprendizaje rápido y sencillo para el usuario, que en todo caso deberá ser menor de 1 hora. \\ \hline
\centering {\bf Importancia} & Vital \\ \hline
\centering {\bf Urgencia} & Inmediatamente \\ \hline
\end{tabular}
\caption{Requisito no funcional 06. Tiempo de aprendizaje por el usuario $<$ 1h.}
\label{enlaceNFR6}
\end{center}
\end{table}

\begin{table}[H]
\begin{center}
\begin{tabular}{|p{3cm}|p{10cm}|} \hline
\centering {\bf NFR-07} & Proporcionar mensajes de apoyo al usuario  \\ \hline\hline
\centering {\bf Versión} & 1.0 (05/06/2022) \\ \hline
\centering {\bf Descripción} & El sistema deberá proporcionar mensajes de apoyo al usuario en todas las operaciones, con el objetivo de facilitar la navegación en la aplicación. \\ \hline
\centering {\bf Importancia} & Quedaría bien \\ \hline
\centering {\bf Urgencia} & Puede esperar \\ \hline
\end{tabular}
\caption{Requisito no funcional 07. Proporcionar mensajes de apoyo al usuario.}
\label{enlaceNFR7}
\end{center}
\end{table}

\begin{table}[H]
\begin{center}
\begin{tabular}{|p{3cm}|p{10cm}|} \hline
\centering {\bf NFR-08} & Uso de tecnologías web y servicios RESTful  \\ \hline\hline
\centering {\bf Versión} & 1.0 (05/06/2022) \\ \hline
\centering {\bf Descripción} & El sistema deberá emplear tecnologías web para realizar la aplicación, incluyendo HTML5, Tailwind como framework de CSS y React.js como framework de JavaScript. Además, todos los servicios serán RESTful, y serán implementados con Spring Boot. \\ \hline
\centering {\bf Importancia} & Vital \\ \hline
\centering {\bf Urgencia} & Inmediatamente \\ \hline
\end{tabular}
\caption{Requisito no funcional 08. Uso de tecnologías web y servicios RESTful.}
\label{enlaceNFR8}
\end{center}
\end{table}

\begin{table}[H]
\begin{center}
\begin{tabular}{|p{3cm}|p{10cm}|} \hline
\centering {\bf NFR-09} & Base de datos NoSQL  \\ \hline\hline
\centering {\bf Versión} & 1.0 (05/06/2022) \\ \hline
\centering {\bf Descripción} & El sistema deberá almacenar todos los datos en una base de datos NoSQL de MongoDB. \\ \hline
\centering {\bf Importancia} & Vital \\ \hline
\centering {\bf Urgencia} & Inmediatamente \\ \hline
\end{tabular}
\caption{Requisito no funcional 09. Base de datos NoSQL.}
\label{enlaceNFR9}
\end{center}
\end{table}

\begin{table}[H]
\begin{center}
\begin{tabular}{|p{3cm}|p{10cm}|} \hline
\centering {\bf NFR-10} & Generación de tokens JWT  \\ \hline\hline
\centering {\bf Versión} & 1.0 (05/06/2022) \\ \hline
\centering {\bf Descripción} & El sistema deberá generar tokens JWT cada vez que un usuario inicie sesión en la aplicación. De esta forma, los datos del usuario no quedan tan accesibles en la web, y será necesario utilizar un algoritmo  {\it hash} para desencriptarlos. En todo caso la contraseña no será almacenada en el token JWT. \\ \hline
\centering {\bf Importancia} & Vital \\ \hline
\centering {\bf Urgencia} & Inmediatamente \\ \hline
\end{tabular}
\caption{Requisito no funcional 10. Generación de tokens JWT.}
\label{enlaceNFR10}
\end{center}
\end{table}
\section{Requisitos de información}

Por último, los requisitos de información \cite{requisitos}, son aquellos requisitos que indican algún tipo de información, la cual es necesaria para el correcto funcionamiento del sistema.
\\

En el caso de esta aplicación, se han localizado dos: Usuario y Documento.
\\

Los requisitos de información pueden ser consultados en el \hyperref[APRequisitosInformacion]{Apéndice C.3. Requisitos de información}.
\section{Casos de uso}

Un caso de uso \cite{casodeuso} es cualquier acción o actividad que se realiza en el sistema. Representan diferentes operaciones que se realizan en la aplicación.
\\

Los casos de uso pueden ser consultados en el \hyperref[APCasosUso]{Apéndice C.4. Casos de uso}.
\cleardoublepage
\chapter{Especificación de requisitos}
\label{enlaceespecificacion}

Para proceder a especificar los requisitos se ha realizado un profundo análisis de los objetivos de la aplicación. Estos requisitos se han buscado teniendo en cuenta aspectos funcionales, de rendimiento, de seguridad, de información, etc.
\\

Las plantillas empleadas para las tablas de los requisitos y casos de uso se encuentran en los apéndices: la \hyperref[enlaceFRQX]{E.1. Plantilla de requisitos funcionales}, la \hyperref[enlaceNFRX]{E.2. Plantilla de requisitos no funcionales}, la \hyperref[enlaceIRQX]{E.3. Plantilla de requisitos de información} y la \hyperref[enlaceUCX]{E.4. Plantilla de casos de uso}.
Estas plantillas están inspiradas en las tablas empleadas en el REM \cite{rem} para la gestión de requisitos.
\\

La importancia de los requisitos será medida según el grado de la necesidad del requisito en la aplicación, siendo en el caso de este proyecto ``Vital" para los requisitos de mayor importancia, y, ``Quedaría bien" para los menos importantes. Por su parte, la urgencia, indicará aquellas partes que se han de implementar con mayor rapidez, siendo ``Inmediatamente" las que tienen una mayor prioridad, y, ``Puede esperar", para los requisitos que no son estrictamente necesarios.

\section{Requisitos funcionales}

Los requisitos funcionales \cite{requisitos} son aquellos que definen una funcionalidad del sistema de software o sus componentes, tanto para las entradas de datos como para el comportamiento en situaciones específicas.
\\

En el caso de esta aplicación se han localizado dos grupos de requisitos funcionales: los asociados con operaciones de usuarios, y los que se relacionan con operaciones asociadas a leyes.
\\

Los requisitos funcionales pueden ser consultados en el \hyperref[APRequisitosFuncionales]{Apéndice C.1. Requisitos funcionales}.
\section{Requisitos no funcionales}
\label{APRequisitosNoFuncionales}

\begin{table}[H]
\begin{center}
\begin{tabular}{|p{3cm}|p{10cm}|} \hline
\centering {\bf NFR-01} & Tiempo de respuesta $<$ 2 segundos  \\ \hline\hline
\centering {\bf Versión} & 1.0 (05/06/2022) \\ \hline
\centering {\bf Descripción} & El sistema deberá responder en menos de 2 segundos, sin tener en cuenta el tiempo que tarde la red en enviar la información de vuelta al usuario. \\ \hline
\centering {\bf Importancia} & Vital \\ \hline
\centering {\bf Urgencia} & Inmediatamente \\ \hline
\end{tabular}
\caption{Requisito no funcional 01. Tiempo de respuesta $<$ 2 segundos.}
\label{enlaceNFR1}
\end{center}
\end{table}

\begin{table}[H]
\begin{center}
\begin{tabular}{|p{3cm}|p{10cm}|} \hline
\centering {\bf NFR-02} & Los datos sensibles del usuario se almacenan encriptados  \\ \hline\hline
\centering {\bf Versión} & 1.0 (05/06/2022) \\ \hline
\centering {\bf Descripción} & El sistema deberá encriptar los datos sensibles del usuario como las contraseñas en la base de datos de lex.gal. Se utilizará un algoritmo SHA-256 para encriptar. \\ \hline
\centering {\bf Importancia} & Vital \\ \hline
\centering {\bf Urgencia} & Inmediatamente \\ \hline
\end{tabular}
\caption{Requisito no funcional 02. Los datos sensibles del usuario se almacenan encriptados.}
\label{enlaceNFR2}
\end{center}
\end{table}

\begin{table}[H]
\begin{center}
\begin{tabular}{|p{3cm}|p{10cm}|} \hline
\centering {\bf NFR-03} & No permitir que un rol realice operaciones de otro  \\ \hline\hline
\centering {\bf Versión} & 1.0 (05/06/2022) \\ \hline
\centering {\bf Descripción} & El sistema deberá restringir aquellas operaciones exclusivas del administrador para el resto de usuarios. \\ \hline
\centering {\bf Importancia} & Vital \\ \hline
\centering {\bf Urgencia} & Inmediatamente \\ \hline
\end{tabular}
\caption{Requisito no funcional 03. No permitir que un rol realice operaciones de otro.}
\label{enlaceNFR3}
\end{center}
\end{table}

\begin{table}[H]
\begin{center}
\begin{tabular}{|p{3cm}|p{10cm}|} \hline
\centering {\bf NFR-04} & Permanecer operativo durante el horario de trabajo \\ \hline\hline
\centering {\bf Versión} & 1.0 (05/06/2022) \\ \hline
\centering {\bf Descripción} & El sistema deberá permanecer operativo durante todo el horario de trabajo de los empleados de lex.gal.  \\ \hline
\centering {\bf Importancia} & Vital \\ \hline
\centering {\bf Urgencia} & Inmediatamente \\ \hline
\end{tabular}
\caption{Requisito no funcional 04. Permanecer operativo durante el horario de trabajo.}
\label{enlaceNFR4}
\end{center}
\end{table}

\begin{table}[H]
\begin{center}
\begin{tabular}{|p{3cm}|p{10cm}|} \hline
\centering {\bf NFR-05} & Tratamiento de documentos en lenguajes de etiquetas  \\ \hline\hline
\centering {\bf Versión} & 1.0 (05/06/2022) \\ \hline
\centering {\bf Descripción} & El sistema deberá restringir el tratamiento de documentos a documentos escritos en lenguajes de etiquetas, como HTML o XML. \\ \hline
\centering {\bf Importancia} & Vital \\ \hline
\centering {\bf Urgencia} & Inmediatamente \\ \hline
\end{tabular}
\caption{Requisito no funcional 05. Tratamiento de documentos en lenguajes de etiquetas.}
\label{enlaceNFR5}
\end{center}
\end{table}

\begin{table}[H]
\begin{center}
\begin{tabular}{|p{3cm}|p{10cm}|} \hline
\centering {\bf NFR-06} & Tiempo de aprendizaje por el usuario $<$ 1h  \\ \hline\hline
\centering {\bf Versión} & 1.0 (05/06/2022) \\ \hline
\centering {\bf Descripción} & El sistema deberá tener un proceso de aprendizaje rápido y sencillo para el usuario, que en todo caso deberá ser menor de 1 hora. \\ \hline
\centering {\bf Importancia} & Vital \\ \hline
\centering {\bf Urgencia} & Inmediatamente \\ \hline
\end{tabular}
\caption{Requisito no funcional 06. Tiempo de aprendizaje por el usuario $<$ 1h.}
\label{enlaceNFR6}
\end{center}
\end{table}

\begin{table}[H]
\begin{center}
\begin{tabular}{|p{3cm}|p{10cm}|} \hline
\centering {\bf NFR-07} & Proporcionar mensajes de apoyo al usuario  \\ \hline\hline
\centering {\bf Versión} & 1.0 (05/06/2022) \\ \hline
\centering {\bf Descripción} & El sistema deberá proporcionar mensajes de apoyo al usuario en todas las operaciones, con el objetivo de facilitar la navegación en la aplicación. \\ \hline
\centering {\bf Importancia} & Quedaría bien \\ \hline
\centering {\bf Urgencia} & Puede esperar \\ \hline
\end{tabular}
\caption{Requisito no funcional 07. Proporcionar mensajes de apoyo al usuario.}
\label{enlaceNFR7}
\end{center}
\end{table}

\begin{table}[H]
\begin{center}
\begin{tabular}{|p{3cm}|p{10cm}|} \hline
\centering {\bf NFR-08} & Uso de tecnologías web y servicios RESTful  \\ \hline\hline
\centering {\bf Versión} & 1.0 (05/06/2022) \\ \hline
\centering {\bf Descripción} & El sistema deberá emplear tecnologías web para realizar la aplicación, incluyendo HTML5, Tailwind como framework de CSS y React.js como framework de JavaScript. Además, todos los servicios serán RESTful, y serán implementados con Spring Boot. \\ \hline
\centering {\bf Importancia} & Vital \\ \hline
\centering {\bf Urgencia} & Inmediatamente \\ \hline
\end{tabular}
\caption{Requisito no funcional 08. Uso de tecnologías web y servicios RESTful.}
\label{enlaceNFR8}
\end{center}
\end{table}

\begin{table}[H]
\begin{center}
\begin{tabular}{|p{3cm}|p{10cm}|} \hline
\centering {\bf NFR-09} & Base de datos NoSQL  \\ \hline\hline
\centering {\bf Versión} & 1.0 (05/06/2022) \\ \hline
\centering {\bf Descripción} & El sistema deberá almacenar todos los datos en una base de datos NoSQL de MongoDB. \\ \hline
\centering {\bf Importancia} & Vital \\ \hline
\centering {\bf Urgencia} & Inmediatamente \\ \hline
\end{tabular}
\caption{Requisito no funcional 09. Base de datos NoSQL.}
\label{enlaceNFR9}
\end{center}
\end{table}

\begin{table}[H]
\begin{center}
\begin{tabular}{|p{3cm}|p{10cm}|} \hline
\centering {\bf NFR-10} & Generación de tokens JWT  \\ \hline\hline
\centering {\bf Versión} & 1.0 (05/06/2022) \\ \hline
\centering {\bf Descripción} & El sistema deberá generar tokens JWT cada vez que un usuario inicie sesión en la aplicación. De esta forma, los datos del usuario no quedan tan accesibles en la web, y será necesario utilizar un algoritmo  {\it hash} para desencriptarlos. En todo caso la contraseña no será almacenada en el token JWT. \\ \hline
\centering {\bf Importancia} & Vital \\ \hline
\centering {\bf Urgencia} & Inmediatamente \\ \hline
\end{tabular}
\caption{Requisito no funcional 10. Generación de tokens JWT.}
\label{enlaceNFR10}
\end{center}
\end{table}
\section{Requisitos de información}

Por último, los requisitos de información \cite{requisitos}, son aquellos requisitos que indican algún tipo de información, la cual es necesaria para el correcto funcionamiento del sistema.
\\

En el caso de esta aplicación, se han localizado dos: Usuario y Documento.
\\

Los requisitos de información pueden ser consultados en el \hyperref[APRequisitosInformacion]{Apéndice C.3. Requisitos de información}.
\section{Casos de uso}

Un caso de uso \cite{casodeuso} es cualquier acción o actividad que se realiza en el sistema. Representan diferentes operaciones que se realizan en la aplicación.
\\

Los casos de uso pueden ser consultados en el \hyperref[APCasosUso]{Apéndice C.4. Casos de uso}.

\end{document}
